\documentclass[11pt]{article}
\usepackage{tcee}
\addbibresource{3A02.bib}
\tcetema[Los economistas clásicos y Marx]{3.A.2}{Los economistas clásicos y Marx.}
\tcefecha{04/11/2026}
\begin{document}
\tceetitulo

\section*{Introducción}
\minutos{2,5}

\begin{itemize}
\item \textbf{\emph{Enganche}}:
  \begin{itemize}
  \item «Cada uno habla de la feria según le va en ella», dice el refrán. El 5 de mayo de 2018, en el bicentenario de su nacimiento, Tréveris descubrió una estatua de bronce de \autor{Karl Marx} de 5,5 metros con su pedestal, regalo del Gobierno chino \parencite{Spiegel2018Marx}; en 2023, la Universidad de Glasgow, donde \autor{Adam Smith} estudió y enseñó, celebró su tricentenario \parencite{Glasgow2023Smith300}.
  \item Los 2 miraron la misma economía, la de la Revolución industrial, y vieron cosas opuestas: \autor{Smith}, una sociedad en la que el interés de cada uno acaba en riqueza para todos; \autor{Marx}, una en la que el capitalista se queda con parte de lo que produce el obrero.
  \item Detrás está la pregunta de esta exposición: \emph{¿de dónde sale la riqueza de una nación y quién se la queda?}
  \end{itemize}

\item \textbf{\emph{Relevancia}}:
  \begin{itemize}
  \item \emph{Para la teoría}: los clásicos plantean las preguntas que la economía sigue haciéndose: el crecimiento \vertema{3.A.43} \vertema{3.A.44}, el reparto de la renta entre los factores \vertema{3.A.45}, el comercio internacional \vertema{3.B.5} y el dinero \vertema{3.A.35}.
  \item \emph{Para la política económica}: el libre comercio (la derogación de las Leyes de Granos en 1846), el patrón oro y la Ley bancaria de 1844 \vertema{3.B.20}, y el temor a los límites del crecimiento, de \autor{Malthus} al Club de Roma \vertema{3.B.31}.
  \item \emph{Para el temario}: es la teoría que corrigen los neoclásicos \vertema{3.A.3}, la que \autor{Keynes} discute con la ley de \autor{Say} \vertema{3.A.4} y la que \autor{Marx} critica desde dentro.
  \item Además, como toda historia del pensamiento económico, enseña a no quedarse con el último manual y a ver por qué unas ideas llegan en unas épocas y no en otras: son 2 de los 4 motivos de \textcite{Schumpeter1954History}.\footnote{Los otros 2: la historia es fuente de ideas (principios que no fructificaron pueden orientar la investigación de hoy) y muestra cómo se desplazan las fronteras de la ciencia (cada época trae sus problemas y su forma de afrontarlos). La evolución del pensamiento económico resulta de la interacción de varios factores: el contexto económico de cada momento, el pensamiento anterior y los avances de otras disciplinas, como la filosofía, las matemáticas, la física o la biología. Conocerla permite entender las raíces del pensamiento actual, analizar los fenómenos económicos pasados y presentes y valorar los distintos programas de investigación.}
  \end{itemize}

\item \textbf{\emph{Contextualización}}:
  \begin{itemize}
  \item Para seguir la historia de las ideas económicas hay que fijar antes 2 cosas: qué estudia la economía, su objeto \vertema{3.A.1}, y desde cuándo hay una ciencia económica. Tomamos el punto de partida de casi todos los manuales: 1776.
  \item Pensamiento económico lo hubo siempre, de \autor{Aristóteles} a los escolásticos, pero no una ciencia económica. Los escritos anteriores al siglo~XVIII tenían 3 limitaciones:
    \begin{itemize}
    \item \emph{no eran solo económicos}: la riqueza aparecía dentro de la moral, la justicia, la teología o el gobierno del Estado;
    \item \emph{eran más normativos que analíticos}: recetaban lo que debía hacer el príncipe más que explicar cómo funciona la economía;
    \item \emph{eran parciales}: estudiaban un problema suelto (la moneda, la balanza comercial, el impuesto) y no un sistema.
    \end{itemize}
  \item En los 2 siglos anteriores, 3 corrientes preparan el terreno (anexo~\ref{anexo:preclasica}):
    \begin{itemize}
    \item los \emph{mercantilistas} identifican la riqueza con los metales preciosos y piden al Estado que proteja la industria y el comercio exterior;\footnote{\textcite{Chang2002} sostiene que los países hoy ricos se industrializaron con políticas «mercantilistas» (aranceles, subvenciones y protección de la industria naciente) y que, una vez arriba, «retiraron la escalera» al recomendar el libre comercio a los países pobres \vertema{3.B.28}. Es la tesis de \textcite{List1841System} frente a \autor{Smith}, y la de \autor{John Stuart Mill} cuando admite la protección de la industria naciente \vertema{3.B.8}.}
    \item la \emph{Escuela de Salamanca} formula la teoría cuantitativa del dinero con \textcite{Azpilcueta1556Cambios}, antes que \textcite{Bodin1568Reponse};
    \item los \emph{fisiócratas}, con el \emph{Tableau économique} de \textcite{Quesnay1758Tableau}, ven la economía como un circuito que se regula solo y defienden el \emph{laissez faire}: son el precedente inmediato, porque ya piensan en un sistema.
    \end{itemize}
  \item En 1776, \textcite{Smith1776Wealth} los reúne en un sistema: con la \emph{Riqueza de las naciones} empieza la \concepto{escuela clásica}, que llega hasta los \emph{Principios} de \textcite{Mill1848Principles} y domina la disciplina hasta la revolución marginalista de los años 70 \vertema{3.A.3}.
  \item Es la época de la Revolución industrial, de la Revolución francesa y de las guerras con Francia, y la escuela es casi enteramente británica. Sus autores escriben en una economía que crece, pero con una población que aumenta deprisa y una tierra que no aumenta.
  \item El nombre de \emph{clásicos} se lo pone su gran crítico, \autor{Marx}, que llama «economía política clásica» a la que, desde \autor{Petty}, estudia las relaciones reales de producción \parencite[cap.~1, nota~33]{Marx1867Kapital}; \textcite[cap.~1]{Keynes1936} lo extenderá a los neoclásicos \vertema{3.A.4}.
  \end{itemize}

\item \textbf{\emph{Problemática}}:
  \begin{itemize}
  \item ¿Qué tienen en común los clásicos: qué estudian, con qué método y con qué ideas?
  \item ¿De dónde sale el valor de las cosas y cómo se reparte el producto entre salarios, beneficios y rentas?
  \item ¿Puede una economía crecer siempre, o va hacia un estado estacionario? ¿Puede faltar demanda?
  \item ¿Qué toma \autor{Marx} de los clásicos y por qué llega a conclusiones opuestas?
  \end{itemize}

\item \textbf{\emph{Estructura}}:
  \begin{itemize}
  \item La exposición va de la escuela en su conjunto a sus autores, y de estos a su crítico:
  \begin{itemize}
  \item \concepto{I. Caracterización general de la escuela clásica} (\ref{sec:caracterizacion}): quiénes la forman, cuándo, con qué ideas comunes y con qué método;
  \item \concepto{II. Los economistas clásicos y su pensamiento} (\ref{sec:autores-clasicos}): \autor{Smith}, \autor{Say}, \autor{Ricardo}, \autor{Malthus} y \autor{John Stuart Mill}, cada uno con su teoría del valor, del crecimiento y la distribución, de los mercados, el dinero y el comercio, y su política;
  \item \concepto{III. Crítica desde el marxismo: \autor{Karl Marx} (1818-1883)} (\ref{sec:marx}): el materialismo histórico, la teoría del valor y la plusvalía, las leyes del movimiento del capitalismo y su valoración.
  \end{itemize}
  \end{itemize}
\end{itemize}

\begin{esquema}[unbreakable]
\textbf{3.A.2. Los economistas clásicos y Marx}\\[2pt]
Introducción\\
I. Caracterización general de la escuela clásica\\
\hspace*{1.5em}1. Caracterización general de la escuela clásica según \autor{Schumpeter}\\
\hspace*{1.5em}2. Autores y rango histórico\\
\hspace*{1.5em}3. Ideas económicas de la economía política clásica\\
\hspace*{1.5em}4. Método\\
II. Los economistas clásicos y su pensamiento\\
\hspace*{1.5em}1. \autor{Adam Smith} (1723-1790)\\
\hspace*{1.5em}2. \autor{Jean-Baptiste Say} (1767-1832)\\
\hspace*{1.5em}3. \autor{David Ricardo} (1772-1823)\\
\hspace*{1.5em}4. \autor{Thomas Robert Malthus} (1766-1834)\\
\hspace*{1.5em}5. \autor{John Stuart Mill} (1806-1873)\\
III. Crítica desde el marxismo: \autor{Karl Marx} (1818-1883)\\
\hspace*{1.5em}1. Contexto, método y obra\\
\hspace*{1.5em}2. Aportaciones\\
\hspace*{1.5em}3. Valoración y legado de \autor{Marx}\\
Conclusión
\end{esquema}

\begin{notaopositor}
Reparto del cante: introducción, \tiempo{2,5}; I, \tiempo{5,5} (I.1, \tiempo{1,5}; I.2, \tiempo{1,5}; I.3, \tiempo{1,5}; I.4, \tiempo{1}); II, \tiempo{13} (II.1, \tiempo{3,5}; II.2, \tiempo{1}; II.3, \tiempo{4,5}; II.4, \tiempo{2}; II.5, \tiempo{2}); III, \tiempo{8} (III.1, \tiempo{1,5}; III.2, \tiempo{5}; III.3, \tiempo{1,5}), y conclusión, \tiempo{1}. \autor{Marx} no baja de \tiempo{8}: es medio título y el tribunal espera ver su ecuación del valor en la pizarra. El modelo de distribución de \autor{Ricardo} no baja de \tiempo{2}: es el gráfico central del tema. La pizarra se organiza en 3 zonas. A la izquierda, el refrán del enganche y las pegatinas, cada una junto a su concepto: «aprendiz de mucho, maestro de nada» (la división del trabajo), «quien tiene tierra tiene guerra» (la renta de la tierra), «donde comen dos, comen tres» (\autor{Malthus}, que dice lo contrario), «trabajar para el obispo» (la plusvalía) y «cría cuervos y te sacarán los ojos» (las leyes del movimiento del capitalismo). En el centro, el gráfico de cada bloque, que se borra al pasar al siguiente: la línea del tiempo (gráfico~\ref{fig:clasicos-tiempo}), los determinantes de la riqueza (gráfico~\ref{fig:riqueza}), la distribución de \autor{Ricardo} (gráfico~\ref{fig:ricardo}), las progresiones de \autor{Malthus} (gráfico~\ref{fig:malthus}), el materialismo histórico y la jornada laboral (gráficos~\ref{fig:materialismo} y~\ref{fig:jornada}). A la derecha se quedan hasta el final las 2 funciones clave, el reparto del producto de \autor{Ricardo} entre rentas, salarios y beneficios \eqref{eq:reparto-ricardo} y el valor de la mercancía de \autor{Marx} \eqref{eq:valor-marx}, que la conclusión compara. Son de profundización (no se cantan): la caja del contexto histórico, la biografía de cada autor más allá de 2 datos, la teoría monetaria en detalle, la ventaja comparativa en extenso, el problema de la transformación paso a paso y los anexos.
\end{notaopositor}

% ==========================================================================
\section{Caracterización general de la escuela clásica}\label{sec:caracterizacion}
\minutos{5,5}

\begin{itemize}
\item Empezamos por la escuela en su conjunto: quiénes la forman, cuándo, qué ideas comparten y con qué método trabajan.
\item \emph{Conclusión anticipada}: la escuela clásica es casi enteramente británica, tiene un líder indiscutido, \autor{Smith}, y comparte un objeto (el crecimiento de la riqueza y su reparto entre clases), unas pocas leyes (población, rendimientos decrecientes, flujo-especie y mercados), una teoría monetaria (el dinero solo fija los precios) y un método deductivo y verbal; dentro de ella, \autor{Malthus} discrepa en la demanda, y \autor{Marx} usa su aparato para criticarla.
\end{itemize}

% --------------------------------------------------------------------------
\subsection{Caracterización general de la escuela clásica según \autor{Schumpeter}}\label{sec:schumpeter}

\begin{itemize}
\item \autor{Schumpeter} reconoce una escuela de pensamiento por 3 rasgos: un \emph{líder}, unos \emph{seguidores} y un \emph{cuerpo doctrinal} propio, con su objeto y su método \parencite{Schumpeter1954History}. La escuela clásica los tiene, y a diferencia de la neoclásica \vertema{3.A.3}, con un líder indiscutido:
  \begin{itemize}
  \item \emph{Líder}: \autor{Adam Smith}. Las preguntas de la \emph{Riqueza de las naciones} (1776) (la división del trabajo, el valor, el reparto, la acumulación, el comercio y el papel del Estado) son las que se hacen todos los demás, y cada uno parte de lo que dejó el anterior: \autor{Ricardo} corrige las incoherencias de \autor{Smith}, y \autor{John Stuart Mill} suaviza las conclusiones de \autor{Ricardo}.
  \item \emph{Seguidores}: \autor{Jean-Baptiste Say} en Francia; \autor{David Ricardo}, \autor{James Mill}, \autor{Robert Torrens}, \autor{Nassau Senior} y \autor{John Stuart Mill} en Gran Bretaña.
    \begin{itemize}
    \item \autor{Thomas Robert Malthus} es clásico en la población y en la renta de la tierra, pero crítico en la demanda: niega que la oferta cree siempre su propia demanda.
    \item \autor{Karl Marx} no es un clásico más, sino su crítico: usa el aparato de la economía política clásica (la teoría del valor trabajo de \autor{Ricardo}, el excedente, las clases y la acumulación) para llegar a conclusiones opuestas a las suyas \parencite{LandrethColander2006Historia} (\ref{sec:marx-contexto}).
    \end{itemize}
  \item \emph{Cuerpo doctrinal}:
    \begin{itemize}
    \item \emph{objeto}: las causas de la riqueza de las naciones, su crecimiento a largo plazo y su reparto entre las 3 clases sociales (trabajadores, capitalistas y terratenientes) en salarios, beneficios y rentas. Es un objeto dinámico y agregado: no hay aún separación entre micro y macroeconomía, y el sujeto no es el individuo, sino la clase. Para \autor{Ricardo}, «determinar las leyes que regulan esta distribución es el problema principal de la economía política» \parencite[prefacio; traducción propia]{Ricardo1817Principles};
    \item \emph{método}: la \emph{economía política}, verbal y sobre todo deductiva, con las instituciones dentro del análisis (\ref{sec:metodo}).
    \end{itemize}
  \end{itemize}
\end{itemize}

% --------------------------------------------------------------------------
\subsection{Autores y rango histórico}\label{sec:autores}

\begin{itemize}
\item \emph{Periodo}: de la \emph{Riqueza de las naciones} \parencite{Smith1776Wealth} a los \emph{Principios} de \textcite{Mill1848Principles}, que fueron el manual de referencia hasta los de \autor{Marshall} (1890) \vertema{3.A.3}. Unos 70 años (casi 100 hasta la revolución marginalista) de pensamiento casi enteramente británico, con \autor{Say} como gran excepción francesa.
\item \emph{Antes de 1776}: los últimos mercantilistas y los fisiócratas ya veían que las piezas de la economía dependen unas de otras, pero nadie había juntado sus aportaciones en un sistema coherente. Lo hace un profesor escocés de filosofía moral, \autor{Adam Smith}.
\item \emph{Los 3 grandes tratados}:
  \begin{itemize}
  \item \emph{An Inquiry into the Nature and Causes of the Wealth of Nations} (1776), de \autor{Adam Smith} (1723-1790) \parencite{Smith1776Wealth};
  \item \emph{On the Principles of Political Economy, and Taxation} (1817), de \autor{David Ricardo} (1772-1823) \parencite{Ricardo1817Principles};
  \item \emph{Principles of Political Economy} (1848), de \autor{John Stuart Mill} (1806-1873) \parencite{Mill1848Principles}.
  \end{itemize}
\item \emph{2 críticos}: \autor{Thomas Robert Malthus} (1766-1834) es clásico en parte (la población, la renta), pero crítico en la demanda; \autor{Karl Marx} (1818-1883) no es un clásico, sino su crítico desde dentro del aparato clásico:
  \begin{itemize}
  \item \autor{Malthus} aporta la teoría de la población, que todos aceptan, pero se aparta de la ortodoxia en la demanda agregada y en la defensa de los terratenientes;
  \item \autor{Marx} toma de los clásicos el valor trabajo, el excedente y las clases, añade conceptos propios y llega a conclusiones opuestas a las suyas.
  \end{itemize}
\item \emph{El final}: hacia 1870 la escuela pierde sus pilares: la doctrina de la población, la teoría del valor trabajo y la del fondo de salarios, de la que \autor{John Stuart Mill} se retracta en 1869 \parencite{Mill1869Thornton}. La revolución marginalista \parencite{Jevons1871Theory} abre otra etapa \vertema{3.A.3}.
\item El gráfico~\ref{fig:clasicos-tiempo} reúne las fechas y las influencias: arriba, las obras y los sucesos que las rodean; abajo, quién influye en quién.
\end{itemize}
\begin{figure}[H]
\centering
\caption{La economía clásica, \autor{Malthus} y \autor{Marx}: cronología e influencias}\label{fig:clasicos-tiempo}
\grafico[\textwidth]{clasicos-tiempo}
\fuente{elaboración propia. Las influencias, a partir de \textcite{LandrethColander2006Historia}, con \autor{Say}, \autor{James Mill} y \autor{Cantillon} añadidos; las flechas indican el sentido de la influencia. En gris, los precursores; en morado, los clásicos; en rojo, \autor{Marx}; en azul, los sucesos.}
\end{figure}

\begin{anotaciones}
\textbf{Contexto histórico: de \autor{Smith} a \autor{Ricardo}} \parencite{HET2026Ricardian}
\begin{itemize}
\item \emph{1776}: el mismo año de la \emph{Riqueza} muere \autor{Hume} y cae \autor{Turgot}, el ministro reformista de Luis~XVI \parencite{HET2026Turgot}. \autor{Smith} escribe para una Gran Bretaña todavía agraria, antes de que la Revolución industrial se note en las cifras.
\item \emph{1789-1815}: la Revolución francesa (en 1793 Luis~XVI es guillotinado) y más de 20 años de guerra con Francia cambian las preguntas: la deuda, los impuestos, el precio del grano y el papel moneda.
\item \emph{La sucesión de \autor{Smith}}: 3 continuadores con proyectos distintos.
  \begin{itemize}
  \item \autor{Say} (1803) ordena la \emph{Riqueza} y la lleva hacia la utilidad y el empresario;
  \item \autor{Malthus} (1798, 1820) se fija en la población y en la demanda; en 1805 es el primer profesor de economía política de Inglaterra, en el colegio de la Compañía de las Indias Orientales de Haileybury \parencite{HET2026Malthus};
  \item \autor{Ricardo} (1817) la depura con un modelo deductivo del reparto, que se convierte en el sistema clásico.
  \end{itemize}
\item \emph{1810-1811}: en la controversia bullionista (\ref{sec:dinero}), \autor{Malthus} y \autor{Ricardo} traban una amistad que dura hasta la muerte de \autor{Ricardo}, en 1823, con discrepancias en casi todo \parencite{Dorfman1989Malthus}.
\end{itemize}
\end{anotaciones}

% --------------------------------------------------------------------------
\subsection{Ideas económicas de la economía política clásica}\label{sec:ideas}

\begin{itemize}
\item Una vez situados los autores en el tiempo, pasamos a lo que comparten: 3 rasgos, 4 leyes y una teoría del dinero.
\end{itemize}

\subsubsection{Rasgos comunes}\label{sec:rasgos}

\begin{itemize}
\item Los clásicos comparten 3 rasgos que los separan de los mercantilistas que los preceden y de los neoclásicos que los siguen \parencite{LandrethColander2006Historia}:
  \begin{itemize}
  \item \emph{Un sistema que tiende a la armonía}: con competencia y libertad, el interés propio lleva al bien común sin que nadie lo planee. Viene del orden natural de los fisiócratas y es la mano invisible de \autor{Smith} (\ref{sec:smith}) \vertema{3.A.22}.
    \begin{itemize}
    \item Frente a la Iglesia de los escolásticos, que juzgaba la moralidad de los tratos, y al mercantilista que pedía la intervención del Estado, los clásicos defienden los mercados libres y la libertad individual: la económica, además de buena en sí, hace que la economía funcione con la máxima eficiencia.
    \item Pero no ignoran los conflictos: entre los terratenientes y quienes ganan con el crecimiento (\autor{Ricardo}), y, a largo plazo, una tendencia al estado estacionario y a los salarios de subsistencia que dio a la economía fama de ciencia sombría.\footnote{La etiqueta de «ciencia lúgubre» (\emph{dismal science}) es de \textcite{Carlyle1849Occasional}, y no por \autor{Malthus}: \autor{Carlyle} reprocha a los economistas que defiendan que todos los hombres son iguales y libres, en un artículo a favor de la esclavitud.}
    \item En esa armonía está el germen de la ortodoxia posterior (los mercados que se vacían) y en sus conflictos, el de la heterodoxia, que llevará \autor{Marx} hasta el final.
    \end{itemize}
  \item \emph{El crecimiento a largo plazo}: les interesa la acumulación y el reparto entre clases, no la asignación de unos recursos dados ni el ciclo.
    \begin{itemize}
    \item Estudian los precios relativos y los mercados para entender cómo cambia el reparto de la renta con el tiempo, y miran también los factores culturales, políticos e históricos del crecimiento.
    \item Suponen que la economía usa todos sus recursos, así que el paro masivo no es su problema: con precios y salarios flexibles los mercados se vacían y el dinero solo determina el nivel de precios \vertema{3.A.4}.
    \item Por eso la teoría clásica, a diferencia de la neoclásica, es \emph{dinámica}: los neoclásicos estudiarán desde los años 70 la asignación de recursos escasos con estática comparativa. Las diferencias (la pregunta, el sujeto, el valor, la distribución, el método y la política) se resumen en la tabla de rasgos de los neoclásicos \vertema{3.A.3}.
    \end{itemize}
  \item \emph{Desconfianza del gobernante}: el mercantilista se veía como el médico de la economía, con un remedio para cada mal, casi siempre una intervención del Estado. \autor{Smith} duda de que el político sepa más que el mercado y de que tenga pericia para sustituir sus valoraciones por las propias. De ahí el \emph{laissez-faire}, con excepciones que cada autor amplía (\autor{Mill}, más que nadie).
  \end{itemize}
\end{itemize}

\subsubsection{Las 4 leyes de los clásicos}\label{sec:leyes}

\begin{itemize}
\item Los clásicos describen el orden natural con 4 leyes, que aceptan en mayor o menor medida:
  \begin{itemize}
  \item \emph{Ley de los rendimientos decrecientes de la tierra}: con la tierra fija, cada nueva dosis de trabajo y capital aplicada a una tierra dada rinde menos que la anterior (\concepto{margen intensivo}), y al crecer la población hay que cultivar tierras cada vez menos fértiles (\concepto{margen extensivo}). La había formulado \textcite{Turgot1767Observations}, y \autor{West}, \autor{Malthus}, \autor{Ricardo} y \autor{Torrens} la redescubren a la vez en 1815 \parencite{Humphrey1997Algebraic} \vertema{3.A.11}.
  \item \emph{Ley de la población}: la población crece mientras el salario supera el de subsistencia, y los alimentos no la siguen \parencite{Malthus1798Essay} (\ref{sec:malthus}).
  \item \emph{Mecanismo flujo-especie}: con patrón oro, un superávit comercial trae oro, sube los precios y acaba con el superávit \parencite{Hume1752Balance} \vertema{3.B.12}.
  \item \emph{Ley de los mercados o ley de \autor{Say}}: no puede haber un exceso de oferta general de bienes, porque los productos se compran con productos (\ref{sec:say}).
  \end{itemize}
\item A ellas se suma la doctrina del \concepto{fondo de salarios}: a corto plazo, el salario es el capital que los capitalistas adelantan a los trabajadores dividido por su número; a largo plazo, la ley de la población lo lleva al de subsistencia (\ref{sec:smith}).
\item Juntas dan la visión clásica del largo plazo: la acumulación hace crecer el empleo, la población lo sigue, la tierra peor encarece el grano y el beneficio cae hasta el estado estacionario; el dinero solo fija los precios.
\end{itemize}

\subsubsection{La teoría monetaria clásica}\label{sec:dinero}

\begin{itemize}
\item \emph{La idea común}: el dinero no es riqueza, sino el medio que la hace circular; a largo plazo, su cantidad solo fija el nivel de precios (la teoría cuantitativa) \vertema{3.A.35}. Es el ataque al mercantilismo: acumular oro no enriquece a un país.
\item \emph{\autor{Hume} (1752): el flujo-especie}. Si un país vende más de lo que compra, le entra oro, suben sus precios frente a los del resto y sus exportaciones se encarecen, hasta que el superávit desaparece \parencite{Hume1752Money,Hume1752Balance}. La balanza comercial se ajusta sola \vertema{3.B.12}.
\item \emph{La controversia bullionista (1797-1821)}: en 1797, para financiar la guerra con Francia, el Banco de Inglaterra deja de convertir sus billetes en oro, y el precio del oro y el tipo de cambio suben \parencite{HET2026Bullion}.
  \begin{itemize}
  \item Los \emph{bullionistas} (\autor{Thornton}, \autor{Horner}, \autor{Ricardo}) lo atribuyen a un exceso de billetes: la prima del oro mide la depreciación del papel \parencite{Thornton1802Enquiry,Ricardo1810HighPrice}. El \emph{Bullion Report} del 8 de junio de 1810 les da la razón y pide volver a la convertibilidad \parencite{BullionReport1810}.
  \item Los \emph{antibullionistas}, con los directores del Banco, lo atribuyen a la guerra (pagos al exterior, malas cosechas): mientras el banco descuente solo letras de operaciones reales, no puede emitir de más (la \concepto{doctrina de las letras reales}).
  \item La ley de 1819, que impulsa \autor{Robert Peel}, devuelve la convertibilidad en 1821.
  \end{itemize}
\item \emph{Currency school frente a banking school (años 40)}: tras las crisis de 1825, 1836 y 1839, el debate es cómo evitar que el banco emita de más.
  \begin{itemize}
  \item La \concepto{\emph{currency school}} (\autor{Overstone}, \autor{Torrens}, \autor{Norman}): los billetes deben variar exactamente como lo haría el oro, uno a uno con las reservas, por encima de una emisión fiduciaria fija. Hereda el bullionismo de \autor{Ricardo}, su precursor: había propuesto pagar los billetes en lingotes \parencite{Ricardo1816Proposals} y separar la emisión de la banca en un banco nacional \parencite{Ricardo1824Plan}, pero murió en 1823, antes del debate.
  \item La \concepto{\emph{banking school}} (\autor{Tooke}, \autor{Fullarton}): con convertibilidad, la cantidad de billetes la decide el comercio; el exceso vuelve al banco (la «ley del reflujo»), y los depósitos y las letras también son dinero \parencite{Tooke1844Inquiry,Fullarton1844Regulation}.
  \item Gana la primera con la \concepto{Ley de \autor{Peel} de 1844} (\emph{Bank Charter Act}), que separa en el Banco de Inglaterra un departamento de emisión y otro de banca. Pero hubo que suspenderla en las crisis de 1847, 1857 y 1866, como predecía la \emph{banking school} \vertema{3.B.20}.
  \end{itemize}
\end{itemize}

% --------------------------------------------------------------------------
\subsection{Método}\label{sec:metodo}

\begin{itemize}
\item Con las ideas comunes a la vista, queda cómo llegan a ellas: el método.
\item \emph{¿Qué es la economía política?} Para los clásicos, una ciencia moral (\autor{Smith} era catedrático de filosofía moral): estudia cómo se producen y se reparten las riquezas y aconseja al legislador. El adjetivo «política» dice que la economía es asunto del Estado y que las instituciones están dentro del análisis \vertema{3.A.1}.
\item \emph{El ensayo}: escriben en prosa, con ejemplos numéricos y sin matemáticas; sus «modelos» son razonamientos con supuestos explícitos.
\item \emph{Diferencias entre autores}:
  \begin{itemize}
  \item \autor{Smith}: equilibra la deducción con la inducción, la historia y la descripción, como los filósofos de la Ilustración escocesa \parencite{HET2026Scottish};
  \item \autor{Ricardo} y sus seguidores: deducción pura a partir de unos pocos supuestos. \textcite{Schumpeter1954History} llamó \concepto{vicio ricardiano} a la costumbre de aplicar directamente a la política las conclusiones de modelos muy simplificados;
  \item \autor{Malthus}: deducción y datos; en la 2.ª edición del \emph{Ensayo} (1803) apoya su tesis en la experiencia de muchos países;
  \item \autor{Senior} y \autor{John Stuart Mill} convierten el método deductivo en doctrina:
    \begin{itemize}
    \item \textcite{Senior1836Outline} reduce la economía a 4 proposiciones generales y separa la \emph{ciencia}, que describe, del \emph{arte}, que aconseja;
    \item \autor{Mill}, en un ensayo publicado en 1836 y recogido en 1844, define la economía como una ciencia abstracta que estudia al hombre solo en cuanto desea riqueza, y la ve como un razonamiento \emph{a priori}, verdadero solo en abstracto \parencite{Mill1836Definition}: es la figura que después se llamará \emph{homo oeconomicus} \vertema{3.A.1}.\footnote{La idea es de \autor{Mill}, pero el nombre de «hombre económico» se lo pusieron sus críticos a finales del siglo~XIX, y la forma latina \emph{homo oeconomicus} se generaliza después \parencite{Persky1995Ethology}.}
    \end{itemize}
  \item \autor{Marx}: el materialismo histórico; las leyes económicas no son naturales, sino propias de cada modo de producción (\ref{sec:marx-contexto}).
  \end{itemize}
\item \emph{Positiva y normativa}: la distinción clásica es la de \autor{Senior} entre ciencia y arte. Los términos son posteriores: la división en ciencia positiva, ciencia normativa y arte es de \textcite{KeynesJN1891Scope} \vertema{3.A.1}. Aplicados a los clásicos son un anacronismo útil, siempre que se avise.
\item \emph{Recapitulación}: una escuela con líder, un objeto (la riqueza y su reparto entre clases), 4 leyes, una teoría monetaria y un método deductivo. Con este marco común, pasamos a ver qué aporta cada autor.
\end{itemize}

% ==========================================================================
\section{Los economistas clásicos y su pensamiento}\label{sec:autores-clasicos}
\minutos{13}

\begin{itemize}
\item Una vez visto lo que comparten, pasamos a lo que aporta cada uno: \autor{Smith}, \autor{Say}, \autor{Ricardo}, \autor{Malthus} y \autor{John Stuart Mill}, con el mismo orden: contexto y obra, valor, crecimiento y distribución, mercados y dinero, comercio y política económica.
\item \emph{Conclusión anticipada}: \autor{Smith} pone la riqueza en la división del trabajo y el mercado; \autor{Say} formula la ley de los mercados; \autor{Ricardo} demuestra que la tierra frena la acumulación y defiende el libre comercio; \autor{Malthus} añade la población y la demanda, y \autor{Mill} separa la producción, sujeta a leyes, del reparto, que depende de las instituciones.
\end{itemize}

% --------------------------------------------------------------------------
\subsection{\autor{Adam Smith} (1723-1790)}\label{sec:smith}

\subsubsection{Contexto, método y obra}\label{sec:smith-contexto}

\begin{itemize}
\item \emph{Vida} \parencite{Otteson2026Smith,Glasgow2026Smith}:
  \begin{itemize}
  \item Es bautizado en Kirkcaldy (Fife, Escocia) el 5 de junio de 1723; su padre, interventor de aduanas, había muerto meses antes.\footnote{Se conoce la fecha del bautizo, no la del nacimiento, y es del calendario juliano: Gran Bretaña no adopta el gregoriano hasta 1752 \parencite{UK1750CalendarAct}. Por eso no coincide de verdad con el 5 de junio de 1883 en que nace \autor{Keynes}.} A los 3 años lo roban unos gitanos y su tío lo recupera enseguida.
  \item Entra en la Universidad de Glasgow con 14 años (1737), donde le marca \autor{Francis Hutcheson}, y pasa de 1740 a 1746 en el Balliol College de Oxford con una beca Snell.
  \item Es catedrático en Glasgow de lógica (1751) y de filosofía moral (1752-1764), y amigo de \autor{David Hume}.
  \item A finales de 1763 acepta ser tutor del joven duque de Buccleuch, hijastro de \autor{Charles Townshend}, a cambio de una pensión vitalicia de 300 libras, y viaja con él por Francia y Suiza (1764-1766): en Ginebra conoce a \autor{Voltaire} y en París trata a \autor{Quesnay} y a \autor{Turgot}.
  \item Escribe la \emph{Riqueza} en Kirkcaldy y Londres; se publica el 9 de marzo de 1776. En 1778 es nombrado comisario de aduanas de Escocia, en Edimburgo; vive con su madre hasta que ella muere (1784) y fallece el 17 de julio de 1790.
  \end{itemize}
\item \emph{Obra}:
  \begin{itemize}
  \item \emph{The Theory of Moral Sentiments} \parencite{Smith1759Theory}: la conducta humana nace de los sentimientos, y la \emph{simpatía}, ponerse en el lugar del otro ante un «espectador imparcial», modera el interés propio.
  \item \emph{An Inquiry into the Nature and Causes of the Wealth of Nations} \parencite{Smith1776Wealth}, en 5 libros: el 1.º, sobre la división del trabajo, el valor y el reparto; el 2.º, sobre el capital; el 3.º, sobre la historia del progreso de la riqueza; el 4.º, sobre la crítica del mercantilismo y de los fisiócratas, y el 5.º, sobre las finanzas del Estado.
  \end{itemize}
\item \emph{Por qué es el líder}: tiene pocas ideas del todo nuevas, pero hace la síntesis que la época necesitaba y une en un sistema las piezas de mercantilistas, fisiócratas y filósofos \parencite{Schumpeter1954History}.
\item \emph{Método}: deducción con historia y observación (\ref{sec:metodo}); sus recomendaciones dependen del contexto: es partidario del \emph{laissez-faire}, pero con excepciones.
\end{itemize}

\subsubsection{La mano invisible}\label{sec:smith-mano}

\begin{itemize}
\item La \concepto{mano invisible} aparece solo 3 veces en \autor{Smith}: en su historia de la astronomía, escrita antes de 1758 y publicada en 1795 (la «mano invisible de Júpiter») \parencite{Smith1795Essays}; en los \emph{Sentimientos morales}, para el reparto de lo necesario para vivir \parencite[parte~IV, cap.~I]{Smith1759Theory}, y una sola vez en la \emph{Riqueza}: quien busca su propia ganancia es «conducido por una mano invisible a promover un fin que no entraba en sus intenciones» \parencite[libro~IV, cap.~II; traducción propia]{Smith1776Wealth}. Es una metáfora secundaria en su obra, que la posteridad ha puesto en el centro \parencite{Rothschild1994Invisible}.
\item \emph{El razonamiento}, en 4 pasos:
  \begin{itemize}
  \item cada uno conoce mejor que nadie su situación y busca mejorarla;
  \item en el intercambio, cada uno ofrece al otro lo que a este le interesa: «no es de la benevolencia del carnicero, del cervecero o del panadero de donde esperamos nuestra cena, sino de que miren por su propio interés» \parencite[libro~I, cap.~II; traducción propia]{Smith1776Wealth};
  \item con competencia, el interés propio lleva a producir lo que la sociedad desea al menor coste: el precio de mercado gravita hacia el natural \vertema{3.A.16};
  \item la armonía se rompe con el monopolio y los acuerdos: «la gente del mismo oficio rara vez se reúne, ni siquiera para divertirse, sin que la conversación acabe en una conspiración contra el público o en alguna maquinación para subir los precios» \parencite[libro~I, cap.~X, parte~II; traducción propia]{Smith1776Wealth}.
  \end{itemize}
\item Es el antecedente intuitivo del primer teorema fundamental de la economía del bienestar (1TFEB), que la economía moderna demuestra con supuestos precisos \vertema{3.A.22}; la teoría de los fallos de mercado precisa cuándo deja de funcionar \vertema{3.A.23}.
\end{itemize}

\subsubsection{Naturaleza y causas de la riqueza de las naciones}\label{sec:smith-riqueza}

\begin{itemize}
\item La primera frase de la \emph{Riqueza} ya ataca al mercantilismo: «el trabajo anual de cada nación es el fondo que la provee originalmente de todas las cosas necesarias y convenientes para la vida» \parencite[introducción; traducción propia]{Smith1776Wealth}. La riqueza es el producto del trabajo, no el oro acumulado, y depende de 2 cosas (gráfico~\ref{fig:riqueza}):
  \begin{itemize}
  \item \emph{la productividad del trabajo}, que crece con la \concepto{división del trabajo}. En una pequeña fábrica de alfileres, 10 operarios que se reparten 18 operaciones hacen más de 48.000 alfileres al día, 4.800 por cabeza, cuando uno solo quizá no haría ni 20 \parencite[libro~I, cap.~I]{Smith1776Wealth} (gráfico~\ref{fig:alfileres}). Las 3 causas: la destreza que da la práctica, el tiempo que se ahorra al no cambiar de tarea y las máquinas que se inventan para cada tarea. Lo contrario del refrán «aprendiz de mucho, maestro de nada»;
  \item \emph{la proporción entre trabajo productivo e improductivo}. Productivo es el que se fija en una mercancía vendible y añade valor a la materia que trabaja (el del obrero); improductivo, el que se consume en el momento en que se presta (el del criado, el del soberano, el del músico). No es que no sea útil: es que no deja un producto con el que pagar más trabajo \parencite[libro~II, cap.~III]{Smith1776Wealth}.
  \end{itemize}
\item \emph{Los límites}: la división del trabajo está limitada por la \emph{extensión del mercado} \parencite[libro~I, cap.~III]{Smith1776Wealth}, de ahí el libre comercio, y exige \emph{acumular capital} antes de producir, el fondo con que se pagan los salarios mientras dura el proceso, de ahí la parsimonia (libro~II).
\item La división del trabajo con rendimientos crecientes es el antecedente del crecimiento endógeno \parencite{Young1928Increasing} \vertema{3.A.44}.
\end{itemize}
\begin{figure}[H]
\centering
\caption{Determinantes de la riqueza de una nación según \autor{Smith}}\label{fig:riqueza}
\grafico[0.75\textwidth]{smith-riqueza}
\fuente{elaboración propia a partir de \textcite[introducción, libro~I, caps.~I-III, y libro~II, cap.~III]{Smith1776Wealth} y \textcite{LandrethColander2006Historia}.}
\end{figure}
\begin{figure}[H]
\centering
\caption{La fábrica de alfileres: 18 operaciones, 10 operarios, más de 48.000 alfileres al día}\label{fig:alfileres}
\grafico[0.5\textwidth]{alfileres}
\fuente{lámina II de «Épinglier», dibujo de Goussier y grabado de Defehrt, en el \emph{Recueil de planches} de la \emph{Encyclopédie} de \autor{Diderot} y \autor{d'Alembert} (1762); dominio público \parencite{Commons1762Epinglier}. Cifras: \textcite[libro~I, cap.~I]{Smith1776Wealth}.}
\end{figure}

\subsubsection{Teoría del valor}\label{sec:smith-valor}

\begin{itemize}
\item \emph{Valor de uso y valor de cambio}: el agua es utilísima y apenas se cambia por nada; el diamante casi no tiene uso y se cambia por muchísimo. Es la \concepto{paradoja del agua y los diamantes} \parencite[libro~I, cap.~IV]{Smith1776Wealth}, que resolverán los marginalistas con la utilidad marginal \vertema{3.A.3}.
\item \emph{La medida del valor}: el \concepto{trabajo comandado}, el que una mercancía permite comprar, es para \autor{Smith} «la medida real del valor de cambio» \parencite[libro~I, cap.~V; traducción propia]{Smith1776Wealth}.
\item \emph{La causa del valor} depende del tipo de sociedad:
  \begin{itemize}
  \item \emph{en el estado primitivo}, sin capital acumulado ni tierra apropiada, el producto es todo del trabajador y los bienes se cambian según el trabajo incorporado: si cazar un castor cuesta el doble que cazar un ciervo, un castor vale 2 ciervos (libro~I, cap.~VI);
  \item \emph{en la sociedad avanzada}, el precio paga a los 3 factores: el \concepto{precio natural} es la suma de los salarios, los beneficios y las rentas a sus tasas naturales, y el \concepto{precio de mercado}, el que fijan la oferta y la demanda efectivas, gravita hacia él (libro~I, cap.~VII) \vertema{3.A.16}.
  \end{itemize}
\item \emph{El problema}: si el precio es la suma de salarios, beneficios y rentas, y estos son a su vez precios, la explicación es circular. Y con la renta \autor{Smith} se contradice: en el capítulo~VI es un componente del precio; en el XI, su efecto. \autor{Ricardo} se queda con la 2.ª idea (\ref{sec:ricardo-valor}).
\end{itemize}

\subsubsection{Crecimiento y distribución}\label{sec:smith-distribucion}

\begin{itemize}
\item \emph{Salarios}:
  \begin{itemize}
  \item \emph{a corto plazo, el fondo de salarios}: los paga un fondo que los capitalistas adelantan, parte del capital circulante; el salario medio es el fondo dividido entre los trabajadores;
  \item \emph{a largo plazo, el salario natural}: tiende al nivel que permite mantener a la familia y reponer la población. Pero en un país que crece el salario es más alto: la «recompensa liberal del trabajo» es efecto y síntoma de la riqueza creciente \parencite[libro~I, cap.~VIII]{Smith1776Wealth}. \autor{Smith} no habla de una «ley de hierro de los salarios»: el nombre se lo pondrá \textcite{Lassalle1863Antwortschreiben} a la idea del salario de subsistencia.
  \end{itemize}
\item \emph{Beneficios}: la remuneración de quien adelanta el capital. Incluyen el interés y una prima por el riesgo: la tasa mínima de beneficio tiene que cubrir las pérdidas ocasionales, y el interés que se puede pagar sale del beneficio neto \parencite[libro~I, cap.~IX]{Smith1776Wealth}. Con la acumulación y la competencia entre capitales, la tasa de beneficio tiende a bajar.
\item \emph{Rentas}: el precio por usar la tierra, un precio de monopolio, lo más que el arrendatario puede pagar (libro~I, cap.~XI). Para \autor{Smith} el residuo es la renta; para \autor{Ricardo}, el beneficio.
\end{itemize}

\subsubsection{El ahorro y el dinero}\label{sec:smith-say}

\begin{itemize}
\item \emph{La ley de \autor{Say} antes de \autor{Say}}: lo que se ahorra cada año se consume tan regularmente como lo que se gasta, y casi al mismo tiempo, pero por otras personas: los trabajadores productivos que lo cobran como salario \parencite[libro~II, cap.~III]{Smith1776Wealth}. El ahorro no frena la demanda: la cambia de manos. La \emph{parsimonia}, y no el gasto, hace crecer el capital.
\item \emph{El dinero}: es la «gran rueda de la circulación», que no produce nada. Los billetes convertibles ahorran oro y lo liberan para comprar capital fuera; para no emitir de más, los bancos deben descontar solo letras de operaciones reales a corto plazo (la doctrina de las letras reales, antecedente de la \emph{banking school}) \parencite[libro~II, cap.~II]{Smith1776Wealth} (\ref{sec:dinero}).
\end{itemize}

\subsubsection{Comercio internacional}\label{sec:smith-comercio}

\begin{itemize}
\item Contra el mercantilismo, el comercio no es un juego de suma cero: ganan los 2 países. Su razón es la del padre de familia prudente, que no intenta hacer en casa lo que le costaría más hacer que comprar \parencite[libro~IV, cap.~II]{Smith1776Wealth}. Un país exporta lo que produce con menos trabajo que los demás (\concepto{ventaja absoluta}) \vertema{3.B.5}.
\item \emph{Ventajas}: amplía el mercado y con él la división del trabajo, y da salida a los excedentes que el mercado interior no absorbe.
\item \emph{Excepciones}: la defensa, porque «la defensa es mucho más importante que la opulencia» (por eso aprueba el Acta de Navegación) \parencite[libro~IV, cap.~II; traducción propia]{Smith1776Wealth}, y los aranceles de represalia si sirven para que el otro país baje los suyos.
\item \emph{El límite}: con ventaja absoluta, un país menos productivo en todo no podría exportar nada. Lo resolverá \autor{Ricardo} (\ref{sec:ricardo-comercio}).
\end{itemize}

\subsubsection{Política económica}\label{sec:smith-politica}

\begin{itemize}
\item \emph{El sistema de libertad natural}: el soberano tiene solo 3 deberes \parencite[libro~IV, cap.~IX, y libro~V, cap.~I]{Smith1776Wealth}:
  \begin{itemize}
  \item la defensa frente a otras sociedades;
  \item la justicia, que protege a cada uno de la opresión de los demás;
  \item las obras e instituciones públicas que ningún particular haría porque no le serían rentables, aunque a la sociedad sí (caminos, puertos, la educación del pueblo).
  \end{itemize}
\item \emph{Los impuestos}: 4 máximas, igualdad (según la renta), certeza, comodidad en el pago y economía en la recaudación \parencite[libro~V, cap.~II]{Smith1776Wealth}.
\item \emph{Contra los privilegios}: las compañías privilegiadas, los gremios y los aranceles son lo que el «sistema mercantil» da a los comerciantes a costa de los consumidores.
\end{itemize}

% --------------------------------------------------------------------------
\subsection{Jean-\autor{Baptiste Say} (1767-1832)}\label{sec:say}

\begin{itemize}
\item \emph{El autor}: industrial del algodón, periodista y profesor en París. Su \emph{Traité d'économie politique} \parencite{Say1803Traite} difunde a \autor{Smith} en el continente, con un orden más claro y 2 novedades: el valor nace de la utilidad, no solo del trabajo (precursor de los neoclásicos \vertema{3.A.3}), y el \emph{empresario}, el que organiza la producción, es un agente distinto del capitalista.
\item \emph{La ley de los mercados}: en el capítulo «Des débouchés», \autor{Say} sostiene que un producto, en cuanto se termina, abre una salida a otros productos por todo su valor: quien vende quiere gastar enseguida el dinero que obtiene \parencite[libro~I, cap.~XV]{Say1803Traite}. Los productos se compran con productos; el dinero es solo un medio de cambio.
  \begin{itemize}
  \item \emph{Implicación}: puede sobrar producción de un bien, si se ha producido mal, pero no de todos a la vez: el exceso de oferta de unos va con un exceso de demanda de otros.
  \item \emph{Quién la formula}: en inglés la expone antes \textcite{MillJames1808Commerce}, el padre de John Stuart: el poder de compra de una nación es su producto anual. La frase «la oferta crea su propia demanda» no es de \autor{Say}: es como la resume \textcite[cap.~2]{Keynes1936} para atacarla \vertema{3.A.4}.
  \end{itemize}
\item \emph{En lenguaje moderno}: en la formulación de \textcite{Lange1942Say}, la ley de \autor{Say} afirma que el valor de los excesos de demanda de los bienes, sin contar el dinero, es idénticamente nulo, y por tanto que nunca puede haber un exceso de oferta general de bienes. Con $n$ bienes de precios $p_k$ y excesos de demanda $z_k(\vect{p})$ \vertema{3.A.21}:
  \begin{equation}\label{eq:say}
  \sum_{k=1}^{n}p_k\,z_k(\vect{p})\equiv 0
  \end{equation}
  No debe confundirse con la ley de \autor{Walras}, que es una identidad que incluye todos los mercados, también el del dinero \vertema{3.A.21}; y es la que \autor{Keynes} critica: en una economía monetaria, un exceso de oferta de bienes puede ir acompañado de un exceso de demanda de dinero \vertema{3.A.4}.
  \begin{itemize}
  \item \textcite{BeckerBaumol1952Classical} distinguen la \concepto{identidad de \autor{Say}}, que se cumple a cualquier precio y deja al dinero sin papel, de la \concepto{igualdad de \autor{Say}}, que solo se cumple en equilibrio y admite excesos de oferta temporales: es la de \autor{John Stuart Mill} (\ref{sec:mill}).
  \end{itemize}
\item \emph{El debate del exceso general de oferta} (\emph{general glut}, 1819-1824): \autor{Sismondi} \parencite{Sismondi1819Nouveaux} y \autor{Malthus} \parencite{Malthus1820Principles} sostienen que puede faltar demanda efectiva (\ref{sec:malthus-demanda}); \autor{Say} \parencite{Say1820Lettres}, \autor{Ricardo} y \autor{James Mill} responden que lo ahorrado se invierte y se gasta igual. Ganan los segundos, y la ley queda en la ortodoxia hasta \autor{Keynes} \vertema{3.A.42}.
\end{itemize}

% --------------------------------------------------------------------------
\subsection{\autor{David Ricardo} (1772-1823)}\label{sec:ricardo}

\subsubsection{Contexto, método y obra}\label{sec:ricardo-contexto}

\begin{itemize}
\item \emph{Vida}: hijo de un agente de bolsa judío sefardí llegado de Holanda, hace fortuna en la Bolsa de Londres, se retira a partir de 1814 y es diputado desde 1819 hasta su muerte \parencite{HET2026Ricardo}. Entra en la economía por la controversia bullionista \parencite{Ricardo1810HighPrice} (\ref{sec:dinero}).
\item \emph{Contexto}: las guerras con Francia habían cerrado el comercio y subido el precio del grano y la renta de la tierra. Con la paz, los terratenientes logran las \concepto{Leyes de Granos} (1815), que frenan la importación de trigo. \autor{Robert Peel} las deroga en 1846, en plena hambruna irlandesa (1845-1849) \parencite{GovUK2026Peel} \vertema{3.B.20}.
\item \emph{Obra}: \emph{An Essay on the Influence of a Low Price of Corn on the Profits of Stock} \parencite{Ricardo1815Essay}, el primer esbozo de su modelo, y los \emph{Principles} \parencite{Ricardo1817Principles}, con 3 ediciones en vida (1817, 1819 y 1821); la 3.ª \parencite{Ricardo1821Principles}, que da la numeración de capítulos que se cita, añade el de la maquinaria.
\item \emph{Método}: pasa de la riqueza de \autor{Smith} al reparto: «determinar las leyes que regulan esta distribución es el problema principal de la economía política» (prefacio). Es el primero que razona con modelos: supuestos simples, deducción y conclusiones que se aplican a la política, lo que \autor{Schumpeter} llama el vicio ricardiano (\ref{sec:metodo}).
\end{itemize}

\subsubsection{Teoría del valor}\label{sec:ricardo-valor}

\begin{itemize}
\item \emph{La regla}: «el valor de una mercancía, o la cantidad de cualquier otra mercancía por la que se cambia, depende de la cantidad relativa de trabajo necesaria para producirla, y no de la mayor o menor remuneración que se paga por ese trabajo» \parencite[cap.~1, secc.~1; traducción propia]{Ricardo1821Principles}. Es el \concepto{trabajo incorporado}, directo e indirecto (el de las herramientas), y vale también en la sociedad avanzada.
\item \emph{Las salvedades}:
  \begin{itemize}
  \item «la utilidad no es la medida del valor de cambio, aunque le es absolutamente esencial» (cap.~1; traducción propia);
  \item los bienes que no se pueden reproducir (una estatua, un vino raro) valen por su escasez, pero son una parte mínima del comercio \parencite[cap.~I, secc.~I]{Ricardo1817Principles};
  \item si los bienes usan capital fijo de distinta duración, una subida de salarios cambia los precios relativos aunque no cambie el trabajo (cap.~1, secc.~4-5): de ahí que \textcite{Stigler1958Ricardo} hable de una teoría del valor trabajo «del 93~\%».
  \end{itemize}
\item \emph{La renta no entra en el precio}: el precio del grano lo fija el trabajo necesario en la tierra peor que se cultiva, la que no paga renta: «el grano no es caro porque se pague una renta, sino que se paga una renta porque el grano es caro» \parencite[cap.~2; traducción propia]{Ricardo1821Principles}. Así resuelve la circularidad de \autor{Smith}.
\item \emph{La medida invariable}: \autor{Ricardo} buscó hasta su muerte una mercancía de valor constante con la que medir los cambios del reparto; no la encontró. \textcite{Sraffa1960Production} la construirá con su «mercancía patrón».
\end{itemize}

\subsubsection{Crecimiento y distribución: el modelo de \autor{Ricardo}}\label{sec:ricardo-distribucion}

\paragraph{a) Idea}
\begin{itemize}
\item ¿Qué pasa con el reparto cuando el país acumula capital y crece la población? Suben las rentas, el salario real se queda en el de subsistencia y la tasa de beneficio cae hasta que la acumulación se para. La culpa es de la tierra, no de los capitalistas ni de los trabajadores: «quien tiene tierra tiene guerra».
\item Es el primer modelo de crecimiento y distribución de la historia del pensamiento, escrito contra las Leyes de Granos \parencite{Ricardo1815Essay}.
\end{itemize}

\paragraph{b) Modelo}
\subparagraph{i) Supuestos}
\begin{itemize}
\item \emph{Un solo sector}, la agricultura, y un solo bien, el grano, que es a la vez producto, capital y salario.
\item \emph{3 clases}: los \emph{terratenientes} cobran la renta y la gastan en lujo; los \emph{trabajadores} cobran el salario y lo consumen; los \emph{capitalistas} adelantan los salarios (el capital es el fondo de salarios, $K=\bar w L$), cobran el beneficio, $B$, y lo ahorran e invierten.
\item \emph{Rendimientos decrecientes de la tierra}: cada nueva dosis de trabajo y capital, $L$, aplicada a una tierra dada rinde menos que la anterior (margen intensivo), y hay que cultivar tierras cada vez peores (margen extensivo) (\ref{sec:leyes}) \vertema{3.A.11}.
\item \emph{Salario natural}: la población sigue al salario (\autor{Malthus}): si supera el de subsistencia, $\bar w$, crece, y a largo plazo el salario real es $\bar w$.
\item \emph{El beneficio es el residuo}: lo que queda tras pagar la renta y los salarios.
\end{itemize}

\subparagraph{ii) Desarrollo}
\begin{itemize}
\item Con una producción $Y=F(L)$ en la tierra disponible, el producto medio es $\mathrm{PMe}_L=F(L)/L$ y el marginal $\mathrm{PMg}_L=F'(L)$, menor y decreciente.
\item \emph{El reparto}: la última dosis no paga renta, así que la tierra cobra lo que el producto medio supera al marginal (\concepto{renta diferencial}: la tierra mejor cobra la diferencia entre lo que rinde y lo que rinde la peor que se cultiva); los trabajadores, el salario de subsistencia, y los capitalistas, el resto:\footnote{La renta diferencial la formulan a la vez, en 1815, \textcite{West1815Essay}, \textcite{Malthus1815Rent}, \textcite{Torrens1815Essay} y \textcite{Ricardo1815Essay}.}
  \funcionclave[eq:reparto-ricardo]{Y=\underbrace{(\mathrm{PMe}_L-\mathrm{PMg}_L)\,L}_{\text{rentas}}+\underbrace{\bar w\,L}_{\text{salarios}}+\underbrace{(\mathrm{PMg}_L-\bar w)\,L}_{\text{beneficios, }B}}
\item \emph{La tasa de beneficio}: como el capital son los salarios adelantados, $K=\bar w L$,
  \begin{equation}\label{eq:beneficio-ricardo}
  \frac{B}{K}=\frac{(\mathrm{PMg}_L-\bar w)\,L}{\bar w\,L}=\frac{\mathrm{PMg}_L-\bar w}{\bar w}
  \end{equation}
\item \emph{La dinámica} (gráfico~\ref{fig:ricardo}): mientras $B/K>0$, los capitalistas acumulan, el fondo de salarios crece, el salario sube por encima de $\bar w$ y la población aumenta hasta devolverlo a $\bar w$. Con más $L$ cae $\mathrm{PMg}_L$: sube la renta y bajan los beneficios y la tasa de beneficio.
\item \emph{El estado estacionario}, en $L^*$: cuando $\mathrm{PMg}_L(L^*)=\bar w$, la tasa de beneficio es nula, nadie acumula, la población deja de crecer y todo el excedente es renta.
\end{itemize}
\begin{figure}[H]
\centering
\caption{El modelo de distribución de \autor{Ricardo}: del reparto en $L_0$ al estado estacionario en $L^*$}\label{fig:ricardo}
\grafico[0.7\textwidth]{ricardo-distribucion}
\fuente{elaboración propia a partir de \textcite{Kaldor1956Alternative} y \textcite[cap.~4]{Blaug1997Economic}; \autor{Ricardo} no dibujó ningún gráfico. Con $F(L)=9L-L^2/2$ y $\bar w=3$: en $L_0=3{,}5$, renta 6,1, beneficios 8,75 y salarios 10,5; en $L^*=6$, renta 18, salarios 18 y beneficios 0. Curvas ilustrativas.}
\end{figure}

\subparagraph{iii) Implicaciones}
\begin{itemize}
\item \emph{Los intereses de los terratenientes se oponen a los de las demás clases}: ganan con todo lo que encarece el grano.
\item \emph{Lo que retrasa el estado estacionario}:
  \begin{itemize}
  \item el \emph{progreso técnico} en la agricultura, que sube la curva del producto marginal;
  \item el \emph{libre comercio} del grano: importar trigo barato permite abandonar las tierras peores, baja la renta y sube el beneficio. De ahí su campaña contra las Leyes de Granos.
  \end{itemize}
\end{itemize}

\paragraph{c) Extensiones}
\begin{itemize}
\item \emph{La versión gráfica es posterior}: la reconstruyen con un solo bien \textcite{Kaldor1956Alternative} y \textcite{Pasinetti1960Mathematical}; es la que se dibuja hoy.
\item \emph{La maquinaria}: en la 3.ª edición (1821) añade un capítulo en el que admite que las máquinas pueden dejar sin trabajo a parte de los obreros \parencite[cap.~31]{Ricardo1821Principles} \vertema{3.A.4}.
\item \emph{La acumulación}: en \autor{Ricardo} se frena por la tierra; en \autor{Marx} será el motor del sistema y el origen de las crisis (\ref{sec:marx-aportaciones}) \vertema{3.A.34}.
\end{itemize}

\paragraph{d) Evidencia empírica}
\begin{itemize}
\item El estado estacionario no llegó: en el siglo~XIX el progreso técnico y la importación de grano, sobre todo tras la derogación de 1846, fueron más deprisa que los rendimientos decrecientes, y el salario real subió \vertema{3.B.20}.
\end{itemize}

\paragraph{e) Valoración}
\begin{itemize}
\item \emph{Lo que acierta}: el razonamiento en el margen (la renta como diferencia) es el germen del análisis marginal \vertema{3.A.3}; y la idea de que un factor fijo frena el crecimiento si no hay progreso técnico reaparece en los rendimientos decrecientes del capital que llevan al estado estacionario de \autor{Solow}, donde el crecimiento a largo plazo solo lo sostiene el progreso técnico \vertema{3.A.43}.
\item \emph{Lo que falla}: infravalora el progreso técnico y supone que la población devuelve siempre el salario a la subsistencia. Con un solo bien, además, evita el problema del valor que \autor{Marx} tendrá que afrontar.
\end{itemize}

\subsubsection{La ley de \autor{Say} y el dinero}\label{sec:ricardo-say}

\begin{itemize}
\item \emph{Ley de \autor{Say}}: la acepta sin reservas: «nadie produce si no es para consumir o vender, y nadie vende si no es para comprar otra mercancía» \parencite[cap.~21; traducción propia]{Ricardo1821Principles}. Lo que puede faltar es beneficio, no demanda; frente a \autor{Malthus}, niega que pueda haber un exceso general de oferta (\ref{sec:say}).
\item \emph{Dinero}: es el gran bullionista: la prima del oro prueba que el Banco de Inglaterra ha emitido de más \parencite{Ricardo1810HighPrice}. Propone un patrón en lingotes \parencite{Ricardo1816Proposals} y un banco nacional que separe la emisión del crédito \parencite{Ricardo1824Plan}: es el precursor de la \emph{currency school}, aunque muere en 1823, antes de su debate (\ref{sec:dinero}).
\end{itemize}

\subsubsection{Comercio internacional: la ventaja comparativa}\label{sec:ricardo-comercio}

\begin{itemize}
\item \emph{La pregunta}: ¿puede comerciar con provecho un país que lo produce todo peor que otro? \autor{Ricardo} responde que sí \parencite[cap.~7]{Ricardo1821Principles} \vertema{3.B.5}.
\item \emph{El ejemplo}: Inglaterra (A) y Portugal (B) producen paño ($X$) y vino ($Y$) solo con trabajo; $a_X^A$ es el trabajo (hombres durante un año) que cuesta una unidad de $X$ en A (tabla~\ref{tab:pano-vino}).
\end{itemize}
\begin{table}[H]
\centering
\caption{El ejemplo de \autor{Ricardo}: trabajo por unidad producida}\label{tab:pano-vino}
\begin{tabular}{>{\raggedright\arraybackslash}m{4.4cm}>{\centering\arraybackslash}m{3.2cm}>{\centering\arraybackslash}m{3.2cm}}
\toprule
\rowcolor{tceeazulpalido} & \textbf{Inglaterra (A)} & \textbf{Portugal (B)}\\
\midrule
Paño ($X$) & $a_X^A=100$ & $a_X^B=90$\\
\addlinespace[0.3em]
Vino ($Y$) & $a_Y^A=120$ & $a_Y^B=80$\\
\addlinespace[0.3em]
\rowcolor{tceesalvia}Coste relativo del paño, $a_X/a_Y$ & \textcolor{tceeazul}{\textbf{0,83}} & 1,125\\
\bottomrule
\end{tabular}
\fuente{\textcite[cap.~7]{Ricardo1821Principles}, con la notación del temario. En azul, la ventaja comparativa de Inglaterra.}
\end{table}
\begin{itemize}
\item Portugal tiene ventaja absoluta en los 2 bienes, pero Inglaterra produce el paño con un coste relativo menor (\concepto{ventaja comparativa}). A se especializa en $X$ y B en $Y$, y el comercio beneficia a los 2 si la relación real de intercambio queda entre los 2 costes relativos:
  \begin{equation}\label{eq:ventaja-comparativa}
  \frac{a_X^A}{a_Y^A}=\frac{100}{120}\;<\;\mathrm{RRI}=\frac{p_X}{p_Y}\;<\;\frac{a_X^B}{a_Y^B}=\frac{90}{80}
  \end{equation}
  Con $\mathrm{RRI}=1$, Inglaterra consigue con 100 unidades de trabajo en paño un vino que le costaría 120, y Portugal, con 80 en vino, un paño que le costaría 90.
\item \emph{Lo que deja abierto}: dónde se fija la RRI entre esos límites; lo cerrará \autor{Mill} con la demanda recíproca (\ref{sec:mill}).
\item \emph{La autoría}: antes que \autor{Ricardo} la apuntan un folleto anónimo de 1814 \parencite{GrancayGrancay2015Considerations} y \textcite{Torrens1815Essay}; la exposición que queda es la de \autor{Ricardo} \parencite{Ruffin2002Ricardo}. El desarrollo completo, en el anexo~\ref{anexo:ricardo-comercio}.
\end{itemize}

\subsubsection{Política económica}\label{sec:ricardo-politica}

\begin{itemize}
\item \emph{Libre comercio}: la derogación de las Leyes de Granos es su gran causa (\ref{sec:ricardo-distribucion}); \autor{Malthus}, en cambio, defiende la protección.
\item \emph{Impuestos}: casi todos acaban recayendo sobre los beneficios y frenan la acumulación; los menos dañinos son los que recaen sobre la renta de la tierra. Prefiere pagar la guerra con impuestos y no con deuda.\footnote{En el capítulo~17 de los \emph{Principles} compara pagar una guerra con un impuesto de 20 millones o con un préstamo cuyos intereses se pagan con un impuesto anual: en teoría, un contribuyente previsor ahorraría lo bastante para pagar los impuestos futuros y las 2 vías serían equivalentes; pero \autor{Ricardo} duda de que la gente razone así \parencite[cap.~17]{Ricardo1821Principles}. Más de 150 años después, \textcite{Barro1974} revive la idea con familias altruistas entre generaciones, y \textcite{Buchanan1976Barro} la bautiza como \concepto{equivalencia ricardiana} \vertema{3.A.39}.}
\item \emph{Dinero}: vuelta a la convertibilidad, con un patrón en lingotes (\ref{sec:ricardo-say}).
\end{itemize}

% --------------------------------------------------------------------------
\subsection{\autor{Thomas Robert Malthus} (1766-1834)}\label{sec:malthus}

\subsubsection{Contexto, método y obra}\label{sec:malthus-contexto}

\begin{itemize}
\item \autor{Ricardo} incorporó la teoría de la población de \autor{Malthus} a su modelo; pasamos a su autor, que fue a la vez su amigo y su principal contradictor.
\item \emph{Vida}: estudia en Cambridge y se ordena clérigo anglicano. En 1805 es el primer profesor de economía política de Inglaterra, en el colegio de la Compañía de las Indias Orientales de Haileybury \parencite{HET2026Malthus}. \autor{Keynes} lo llamará «el primero de los economistas de Cambridge» \parencite{Keynes1933Essays}.
\item \emph{Contexto}: desde 1789 Gran Bretaña importa grano de forma regular \parencite{TNA2026Grain}, y la producción agraria por habitante, que había crecido hasta los años 70, empieza a caer \parencite{BoE2017Millennium}. Los utópicos \autor{Godwin} y \autor{Condorcet} creían que la miseria era culpa de las instituciones y se podía abolir: \autor{Malthus} escribe para refutarlos.
\item \emph{Obra}:
  \begin{itemize}
  \item \emph{An Essay on the Principle of Population} \parencite{Malthus1798Essay}, sin firma en 1798. Tiene 6 ediciones en vida de su autor (1798-1826); la de 1803 es casi otro libro, mucho más empírica, y hoy se leen la 1.ª o la 6.ª \parencite{Malthus1826Essay};
  \item \emph{An Inquiry into the Nature and Progress of Rent} \parencite{Malthus1815Rent}, con la renta diferencial;
  \item \emph{Principles of Political Economy} \parencite{Malthus1820Principles}, con la demanda efectiva.
  \end{itemize}
\item \emph{Método}: deducción a partir de 2 postulados, con datos: en 1798, las tablas de \autor{Price} y \autor{Süssmilch} y la duplicación de la población de las colonias americanas en 25 años; en 1803, la experiencia de muchos países.
\end{itemize}

\subsubsection{El principio de la población}\label{sec:malthus-poblacion}

\paragraph{a) Idea}
\begin{itemize}
\item ¿Puede una sociedad acabar con la pobreza cambiando sus instituciones? \autor{Malthus} responde que no: la población tiende a crecer más deprisa que los alimentos, y la diferencia la cierran la miseria y el vicio. Es justo lo contrario del refrán «donde comen dos, comen tres».
\item El propio \autor{Malthus} reconoce que la idea estaba en \autor{Hume}, en \autor{Smith} y en \autor{Robert Wallace} \parencite[cap.~1]{Malthus1798Essay}, y ya \textcite{Cantillon1755Essai} ligaba la población a los medios de vida.
\end{itemize}

\paragraph{b) Modelo}
\subparagraph{i) Supuestos}
\begin{itemize}
\item Los alimentos son necesarios para la existencia del hombre.
\item La pasión entre los sexos es necesaria y seguirá más o menos como está \parencite[cap.~1]{Malthus1798Essay}.
\end{itemize}

\subparagraph{ii) Desarrollo}
\begin{itemize}
\item Sin frenos, la población se duplica cada 25 años, en \emph{progresión geométrica} (1, 2, 4, 8, 16, 32…); los alimentos crecen, como mucho, en \emph{progresión aritmética} (1, 2, 3, 4, 5, 6…). En 225 años, población y alimentos estarían «como 512 a 10» \parencite[cap.~2; traducción propia]{Malthus1798Essay} (gráfico~\ref{fig:malthus}).
\item La brecha la cierran 2 tipos de \concepto{frenos}:
  \begin{itemize}
  \item \emph{positivos}, que suben la mortalidad: hambrunas, epidemias y guerras;
  \item \emph{preventivos}, que bajan la natalidad: en 1798, el vicio; desde 1803, también la \concepto{contención moral}, retrasar el matrimonio por prudencia, sin relaciones fuera de él \parencite[libro~I, cap.~II]{Malthus1826Essay}.
  \end{itemize}
\end{itemize}
\begin{figure}[H]
\centering
\caption{El principio de la población de \autor{Malthus}}\label{fig:malthus}
\grafico[0.6\textwidth]{malthus-poblacion}
\fuente{elaboración propia a partir de \textcite[caps.~1-2]{Malthus1798Essay}. Índices con el periodo inicial igual a 1; cada periodo, 25 años.}
\end{figure}

\subparagraph{iii) Implicaciones}
\begin{itemize}
\item Para el \autor{Malthus} de 1798, la miseria y el vicio mantienen la población al nivel de los alimentos, y esa desigualdad es «insuperable» en el camino hacia la perfección de la sociedad \parencite[caps.~1-2; traducción propia]{Malthus1798Essay}.
\item El salario real no puede subir de forma duradera por encima de la subsistencia: si sube, nacen más hijos y la mayor oferta de trabajo lo devuelve abajo. Es la pieza que \autor{Ricardo} pone en su modelo (\ref{sec:ricardo-distribucion}).
\item Las \emph{leyes de pobres}, que ayudan según el número de hijos, agravan el problema: \autor{Malthus} pide abolirlas poco a poco.
\end{itemize}

\paragraph{c) Extensiones}
\begin{itemize}
\item \autor{Darwin} y \autor{Wallace} tomaron del \emph{Ensayo} la idea de la lucha por la existencia.
\item Los \emph{neomaltusianos} y el Club de Roma recuperan el temor a los límites del crecimiento, ahora por los recursos y el clima \vertema{3.B.31}.
\end{itemize}

\paragraph{d) Evidencia empírica}
\begin{itemize}
\item Describe bien las economías preindustriales: durante siglos, más renta trajo más población y no más renta por habitante (la \emph{trampa maltusiana}) \vertema{3.A.43}.
\item En Inglaterra, entre 1700 y 1850 la población se multiplica por 3,3 y la producción agraria por 3,1: las 2 crecen de forma geométrica, no la producción en progresión aritmética \parencite{BoE2017Millennium,Broadberry2015British} (gráfico~\ref{fig:malthus-datos}).
  \begin{itemize}
  \item Pero la producción agraria interior por habitante cae, de un índice de 125 en los años 70 del siglo~XVIII a 93 en 1850: la diferencia la cubren las importaciones de grano, pagadas con las exportaciones de la industria.
  \item Y el salario real sube en el siglo~XIX sin que los frenos maltusianos lo devuelvan a la subsistencia: el siglo~XIX no confirma a \autor{Malthus}, como recuerda una pregunta oficial del test (anexo~\ref{anexo:test}).
  \end{itemize}
\end{itemize}
\begin{figure}[H]
\centering
\caption{Inglaterra, 1700-1869: población y producción agraria (índice 1700-1709 = 100)}\label{fig:malthus-datos}
\grafico[0.75\textwidth]{malthus-datos}
\fuente{Bank of England, \emph{A millennium of macroeconomic data}, hojas A2 y A3 \parencite{BoE2017Millennium}, con datos de \textcite{Broadberry2015British}.}
\end{figure}

\paragraph{e) Valoración}
\begin{itemize}
\item \emph{El progreso técnico}: los fertilizantes, la maquinaria y las semillas multiplicaron los rendimientos de la tierra. En tiempos de \autor{Malthus} no había una teoría del progreso técnico; hoy la dan los modelos de crecimiento endógeno \vertema{3.A.44}.
\item \emph{La natalidad}: no previó que, al crecer la renta y la educación, las familias tendrían menos hijos por decisión propia: la transición demográfica.
\item \emph{La vuelta del argumento}: \textcite{Kremer1993} y \textcite{Simon1981Ultimate} sostienen lo contrario: más población da más ideas, y las ideas, más crecimiento \vertema{3.A.44}.
\end{itemize}

\subsubsection{La demanda efectiva y la política económica}\label{sec:malthus-demanda}

\begin{itemize}
\item \emph{La tesis}: para que el producto se venda a precios que cubran los costes hace falta una \concepto{demanda efectiva} suficiente; si los capitalistas ahorran demasiado, la inversión crea más capacidad de la que el consumo puede comprar, y aparece un \emph{exceso general de oferta} \parencite{Malthus1820Principles}.
\item \emph{Los trabajadores no bastan}: cobran la subsistencia y no pueden comprar todo lo que producen. Hace falta una clase de \emph{consumidores improductivos} que gasten sin producir: los terratenientes, el Estado, los criados.
\item \emph{La polémica con \autor{Ricardo}}: este responde que lo ahorrado se invierte y se gasta igual (\ref{sec:say}). \autor{Malthus} no tenía una teoría del atesoramiento ni del tipo de interés, y perdió el debate. Discrepan también en las Leyes de Granos: \autor{Malthus} defiende la protección y \autor{Ricardo}, el grano barato \parencite{Ricardo1815Essay}.
\item \emph{El eco en \autor{Keynes}}: no le dedicó la \emph{Teoría general}, pero en su capítulo~23 lo reivindica como precursor de la demanda efectiva y lamenta que se impusiera la línea de \autor{Ricardo} \parencite[cap.~23]{Keynes1936}, como ya había hecho en su ensayo biográfico \parencite{Keynes1933Essays} \vertema{3.A.4}. Hoy se le cuenta entre los precursores de las teorías del subconsumo \vertema{3.A.42}.
\end{itemize}

% --------------------------------------------------------------------------
\subsection{\autor{John Stuart Mill} (1806-1873)}\label{sec:mill}

\subsubsection{Contexto, método y obra}\label{sec:mill-contexto}

\begin{itemize}
\item \emph{Vida}: educado en casa por su padre, \autor{James Mill}, empieza el griego a los 3 años y estudia economía política con \autor{Ricardo} como texto; en otoño de 1826 sufre una crisis que le aparta del utilitarismo estricto de \autor{Bentham} \parencite{Mill1873Autobiography}. Trabaja en la Compañía de las Indias Orientales, se casa con \autor{Harriet Taylor}, que influye en sus ideas sociales, y es diputado.
\item \emph{Obra}: \emph{Essays on Some Unsettled Questions of Political Economy} \parencite{Mill1844Essays}, escritos hacia 1830, con el de método de 1836 (\ref{sec:metodo}); \emph{Principles of Political Economy} \parencite{Mill1848Principles}, el manual de referencia hasta \autor{Marshall}, y \emph{On Liberty} \parencite{Mill1859Liberty}.
\item \emph{Contexto}: las revoluciones de 1848 y el auge del socialismo. \autor{Mill} quiere un liberalismo con reformas que haga innecesaria la revolución.
\item \emph{Lugar en la escuela}: consolida y suaviza a \autor{Ricardo}. Con sus excepciones al \emph{laissez-faire} es un precursor de la economía del bienestar, aunque la expresión «fallo de mercado» es muy posterior \parencite{Bator1958Anatomy} \vertema{3.A.23}.
\end{itemize}

\subsubsection{Valor, crecimiento y distribución}\label{sec:mill-distribucion}

\begin{itemize}
\item \emph{El valor, «completo»}: «felizmente, no queda nada en las leyes del valor que un escritor actual o futuro tenga que aclarar; la teoría está completa» \parencite[libro~III, cap.~I, §~1; traducción propia]{Mill1848Principles}. Distingue 3 casos según la oferta: los bienes de cantidad fija (su precio lo fija la demanda), los que se pueden reproducir a coste constante (lo fija el coste de producción) y los que se producen a coste creciente, como el grano (lo fija el coste en el margen).
\item \emph{Leyes de la producción y leyes de la distribución}: las de la producción «participan del carácter de las verdades físicas»; la distribución de la riqueza, en cambio, «es solo cuestión de instituciones humanas» \parencite[libro~II, cap.~I; traducción propia]{Mill1909Principles}. La sociedad puede cambiar el reparto: es el hueco por el que entran sus reformas.
\item \emph{El fondo de salarios}: en los \emph{Principles} lo mantiene, y de él deduce que los sindicatos no pueden subir el salario de todos. En 1869, al reseñar a \autor{Thornton}, lo abandona: el capitalista puede pagar más salarios a costa de su propio consumo \parencite{Mill1869Thornton}. Con ello cae una pieza de la escuela (\ref{sec:autores}).
\item \emph{El estado estacionario, deseable}: para \autor{Ricardo} era una amenaza; para \autor{Mill}, una oportunidad, si la población se contiene y el reparto mejora: «una condición estacionaria del capital y la población no implica un estado estacionario del progreso humano» \parencite[libro~IV, cap.~VI; traducción propia]{Mill1909Principles}.
\end{itemize}

\subsubsection{La ley de \autor{Say} y el dinero}\label{sec:mill-say}

\begin{itemize}
\item \emph{Ley de \autor{Say}}: en el trueque no puede haber un exceso general de oferta. Con dinero, sí, de forma transitoria: en una crisis todos quieren vender y nadie comprar, porque todos quieren tener dinero; hay un exceso de oferta de bienes y un exceso de demanda de dinero \parencite[ensayo~II]{Mill1844Essays}. Es la \emph{igualdad} de \autor{Say}, no la identidad (\ref{sec:say}); rechaza, en cambio, el exceso general permanente de \autor{Malthus}.
\item \emph{Dinero}: completa el flujo-especie de \textcite{Hume1752Balance} con los tipos de interés: si sale oro, el banco sube el tipo de descuento, entra capital y el ajuste se acorta \vertema{3.B.12}. En la disputa de los años 40 se alinea con la \emph{banking school} \parencite{HET2026Bullion}: concede a la Ley de \autor{Peel} cierto éxito, pero cree que en conjunto agrava las crisis comerciales \parencite[libro~III, cap.~XXIV]{Mill1909Principles} (\ref{sec:dinero}) \vertema{3.A.35}.
\end{itemize}

\subsubsection{Comercio internacional y política económica}\label{sec:mill-comercio}

\begin{itemize}
\item \emph{La demanda recíproca}: \autor{Ricardo} dejaba abierta la relación real de intercambio. \autor{Mill} la cierra con la «ecuación de la demanda internacional»: la RRI es la que iguala el valor de las exportaciones de cada país con el de sus importaciones, y se acerca más a los costes del país cuya demanda del bien del otro es más intensa: gana más quien menos necesita al otro \parencite[ensayo~I]{Mill1844Essays}. \textcite{Marshall1879Pure} lo dibujará con las curvas de oferta recíproca \vertema{3.B.5}.
\item \emph{La industria naciente}: «el único caso en que, por meros principios de economía política, pueden defenderse los aranceles protectores es cuando se imponen temporalmente (sobre todo en una nación joven y en ascenso) con la esperanza de aclimatar una industria extranjera» \parencite[libro~V, cap.~X; traducción propia]{Mill1909Principles} \vertema{3.B.8}.
\item \emph{Funciones del Estado}: distingue las necesarias (defensa, justicia, propiedad) de las facultativas, y admite excepciones al \emph{laissez-faire}: la educación, la protección de los menores y los servicios que ningún particular prestaría (libro~V).
\item \emph{Impuestos}: proporcionales sobre la renta, con un mínimo exento; progresivos sobre las herencias.
\item \emph{Reformas sociales}: el voto de las mujeres, las cooperativas de trabajadores y la educación como freno a la natalidad.
\item \emph{Recapitulación del apartado}: de \autor{Smith} a \autor{Mill}, la escuela pasa de la riqueza al reparto y de la armonía al estado estacionario; \autor{Mill} abre la puerta a reformar el reparto. \autor{Marx} irá más lejos: el reparto capitalista no se reforma, se supera.
\end{itemize}

% ==========================================================================
\section{Crítica desde el marxismo: \autor{Karl Marx} (1818-1883)}\label{sec:marx}
\minutos{8}

\begin{itemize}
\item Vistos los clásicos, pasamos a su gran crítico, que usa su mismo aparato, el valor trabajo de \autor{Ricardo}, para llegar a la conclusión contraria.
\item \emph{Conclusión anticipada}: para \autor{Marx}, el beneficio es trabajo no pagado (\emph{plusvalía}); la competencia obliga a acumular y a mecanizar, lo que hace caer la tasa de ganancia y trae crisis cada vez más graves, hasta el fin del capitalismo. Su teoría del valor no resiste el paso a los precios (el \emph{problema de la transformación}), pero sus preguntas sobre la acumulación, el reparto y las crisis siguen vivas.
\end{itemize}

% --------------------------------------------------------------------------
\subsection{Contexto, método y obra}\label{sec:marx-contexto}

\begin{itemize}
\item \emph{Vida}: nace en Tréveris (Prusia renana), se doctora en filosofía (1841) y dirige un periódico liberal hasta su cierre (1843). Ese año se va a París, donde en 1844 empieza su colaboración con \autor{Friedrich Engels}; expulsado de París (1845), de Bruselas (1848) y de Colonia (1849), desde 1849 vive en Londres, donde estudia a los economistas en la biblioteca del Museo Británico.
\item \emph{3 fuentes}: la dialéctica de \autor{Hegel} (la historia avanza por contradicciones), vuelta del revés con el materialismo de \autor{Feuerbach}; el socialismo francés, y la economía política inglesa, sobre todo \autor{Ricardo}.
\item \emph{Contexto}: la industrialización concentra a los obreros en las ciudades, con jornadas larguísimas, trabajo infantil y crisis periódicas (1825, 1836, 1847, 1857).
\item \emph{Obra económica}:
  \begin{itemize}
  \item el \emph{Manifiesto comunista}, con \autor{Engels} \parencite{MarxEngels1848Manifest};
  \item la \emph{Contribución a la crítica de la economía política} \parencite{Marx1859Kritik}, cuyo prólogo resume el materialismo histórico;
  \item \emph{El capital}: el libro~I, sobre la producción \parencite{Marx1867Kapital}; el libro~2.º, sobre la circulación \parencite{Marx1885Kapital}, y el 3.º, sobre el proceso en su conjunto \parencite{Marx1894Kapital}, los publica \autor{Engels} con sus borradores, y las \emph{Teorías de la plusvalía}, \autor{Kautsky} \parencite{Marx1905Theorien}.\footnote{Los capítulos de \emph{El capital} se citan con la numeración de la edición alemana, la de la traducción castellana de Pedro Scaron (Siglo~XXI); la inglesa de 1887 suma 2 a partir del capítulo~5.}
  \end{itemize}
\item \emph{Método}: el materialismo histórico (\ref{sec:marx-aportaciones}), con una exposición que va de lo abstracto (la mercancía) a lo concreto (los precios, la renta, el crédito). Sus leyes no son naturales ni eternas, como las de los clásicos: son las del modo de producción capitalista, que es histórico y pasará.
\item \emph{¿Clásico?} \autor{Marx} no es un clásico más, sino su crítico: usa el aparato de la economía política clásica (la teoría del valor trabajo de \autor{Ricardo}, el excedente, las clases y la acumulación) para llegar a conclusiones opuestas a las suyas \parencite{LandrethColander2006Historia}. Es él quien les pone el nombre (\ref{sec:schumpeter}), y el subtítulo de \emph{El capital} es «crítica de la economía política». La tabla~\ref{tab:clasicos-marx} resume lo que toma y lo que cambia.
\end{itemize}
\begin{table}[H]
\centering
\caption{Los clásicos y \autor{Marx}: lo que toma y lo que cambia}\label{tab:clasicos-marx}
\small
\begin{tabular}{>{\raggedright\arraybackslash}m{2.4cm}>{\raggedright\arraybackslash}m{5.6cm}>{\raggedright\arraybackslash}m{6cm}}
\toprule
\rowcolor{tceeazulpalido}\textbf{Rasgo} & \textbf{Clásicos} & \textbf{\autor{Marx}}\\
\midrule
Valor & Trabajo incorporado (\autor{Ricardo}) o coste de producción (\autor{Smith}, \autor{Mill}) & Trabajo socialmente necesario; la fuerza de trabajo crea más valor del que cuesta\\
\addlinespace[0.3em]
Clases & 3: trabajadores, capitalistas y terratenientes & \textcolor{tceerojo}{\textbf{2}}: burguesía y proletariado, por la propiedad de los medios de producción\\
\addlinespace[0.3em]
Conflicto & Terratenientes frente a capitalistas, en un sistema que tiende a la armonía & \textcolor{tceerojo}{\textbf{Capital frente a trabajo}}: la explotación\\
\addlinespace[0.3em]
Leyes & Naturales y universales & \textcolor{tceerojo}{\textbf{Históricas}}: propias del capitalismo\\
\addlinespace[0.3em]
Dinámica & Beneficio decreciente y estado estacionario & Tasa de ganancia decreciente, crisis y fin del capitalismo\\
\addlinespace[0.3em]
El capitalista & Motor de la acumulación y del crecimiento & Motor de la acumulación y causa de la miseria\\
\bottomrule
\end{tabular}
\fuente{elaboración propia a partir de \textcite{LandrethColander2006Historia}. En rojo, lo que cambia.}
\end{table}

% --------------------------------------------------------------------------
\subsection{Aportaciones}\label{sec:marx-aportaciones}

\begin{itemize}
\item Las aportaciones de \autor{Marx} van de la sociedad en su conjunto al capitalismo y a su futuro: a) qué mueve la historia; b) de dónde sale el beneficio; c) cómo crece el capital; d) adónde va el capitalismo, y e) qué hacer.
\end{itemize}

\paragraph{a) El materialismo histórico}
\begin{itemize}
\item \emph{La idea}: lo primero que tiene que hacer una sociedad es producir para vivir, y cómo lo hace explica lo demás: no es la conciencia de los hombres la que determina su ser, sino su ser social el que determina su conciencia \parencite[prólogo]{Marx1859Kritik}.
\item \emph{Las piezas} (gráfico~\ref{fig:materialismo}):
  \begin{itemize}
  \item las \concepto{fuerzas productivas}: los medios de producción (tierra, herramientas, máquinas) y la fuerza de trabajo, con sus conocimientos;
  \item las \concepto{relaciones de producción}: quién es dueño de los medios de producción y cómo se reparte el producto. Su conjunto es la \concepto{estructura económica}, «la base real sobre la que se levanta una superestructura jurídica y política y a la que corresponden formas determinadas de conciencia social» \parencite[prólogo; traducción propia]{Marx1859Kritik};
  \item el \concepto{modo de producción}: la combinación de unas fuerzas y unas relaciones de producción. El prólogo enumera el asiático, el antiguo, el feudal y el burgués moderno.
  \end{itemize}
\item \emph{El motor}: cuando las fuerzas productivas crecen y chocan con las relaciones de producción, estas se vuelven una traba y se abre una época de revolución social. Los actores son las clases: «la historia de todas las sociedades hasta nuestros días es la historia de la lucha de clases» \parencite[cap.~I]{MarxEngels1848Manifest}. En el capitalismo, la burguesía, dueña del capital, frente al proletariado, que solo tiene su fuerza de trabajo.
\item \emph{El fetichismo de la mercancía}: en el mercado, una relación social entre personas «adopta… la forma fantasmagórica de una relación entre cosas», los precios de las mercancías, que parecen tener valor por sí mismas \parencite[libro~I, cap.~1, §~4; traducción propia]{Marx1867Kapital}.
\end{itemize}
\begin{figure}[H]
\centering
\caption{El materialismo histórico}\label{fig:materialismo}
\grafico[0.7\textwidth]{materialismo-historico}
\fuente{esquema propio a partir de \textcite[prólogo]{Marx1859Kritik} y \textcite[cap.~I]{MarxEngels1848Manifest}.}
\end{figure}

\paragraph{b) La teoría del valor y de la plusvalía}

\subparagraph{i) Idea}
\begin{itemize}
\item Si en el mercado se cambian equivalentes, ¿de dónde sale el beneficio? De una mercancía especial, la fuerza de trabajo, que crea más valor del que cuesta.
\end{itemize}

\subparagraph{ii) Supuestos}
\begin{itemize}
\item \emph{La mercancía} tiene un \emph{valor de uso} (sirve para algo) y un \emph{valor}, que se expresa en el valor de cambio.
\item \emph{El valor} de una mercancía es el \concepto{tiempo de trabajo socialmente necesario} para producirla, con la técnica normal y la destreza y la intensidad medias: quien tarda más no crea más valor \parencite[libro~I, cap.~1]{Marx1867Kapital}.
\item \emph{El doble carácter del trabajo}: el mismo trabajo es \emph{concreto} (tejer, hacer pan), y crea valores de uso, y \emph{abstracto} (gasto de trabajo humano en general), y crea valor (libro~I, cap.~1, §~2).
\item \emph{Competencia}: en el libro~I, las mercancías se cambian por su valor, y el trabajador es libre y vende su fuerza de trabajo por lo que vale.
\end{itemize}

\subparagraph{iii) Desarrollo}
\begin{itemize}
\item \emph{Fuerza de trabajo y trabajo}: el obrero vende su \concepto{fuerza de trabajo}, no su trabajo. Su valor es el de los bienes que necesita para vivir y reproducirse, con un elemento histórico y moral (libro~I, cap.~4); pero en una jornada produce más valor que ese: la diferencia es la \concepto{plusvalía}, que se apropia el capitalista \parencite[libro~I, caps.~4-7]{Marx1867Kapital}.
  \begin{itemize}
  \item Volvamos a la feria del enganche: en una jornada de 10 horas, con las 5 primeras el obrero repone su salario (\emph{trabajo necesario}) y las otras 5 las trabaja para el capitalista (\emph{plustrabajo}) (gráfico~\ref{fig:jornada}). Es lo que el refranero llama «trabajar para el obispo».
  \end{itemize}
\item \emph{El capital} se divide en 2 partes (libro~I, cap.~6):
  \begin{itemize}
  \item el \concepto{capital constante}, $c$: las máquinas y las materias primas, trabajo pasado (\emph{trabajo muerto}). Transfiere su valor al producto, pero no crea valor nuevo;
  \item el \concepto{capital variable}, $v$: los salarios, el valor de la fuerza de trabajo. El \emph{trabajo vivo} que compra repone su valor y añade la plusvalía, $\mathit{pl}$.
  \end{itemize}
\item \emph{El valor de una mercancía}:\footnote{Notación de \autor{Marx}: $W$ (\emph{Wert}) es el valor de la mercancía; en el resto del temario, $W$ es el salario nominal y $\overline W$ la renta del consumidor. Las primas de $\mathit{pl}'$ y $g'$ son las de \emph{El capital} y evitan confundirlas con el progreso técnico, $g$.}
  \funcionclave[eq:valor-marx]{W=\underbrace{c}_{\text{trabajo muerto}}+\underbrace{v+\mathit{pl}}_{\text{trabajo vivo}}}
\item \emph{Las 3 razones}: la \concepto{tasa de plusvalía} o de explotación, la \concepto{composición orgánica del capital} y la \concepto{tasa de ganancia} (libro~I, cap.~7; libro~III, caps.~1-3):
  \begin{equation}\label{eq:tasas-marx}
  \mathit{pl}'=\frac{\mathit{pl}}{v},\qquad \frac{c}{v},\qquad g'=\frac{\mathit{pl}}{c+v}=\frac{\mathit{pl}'}{1+c/v}
  \end{equation}
  \begin{itemize}
  \item \emph{En el ejemplo}: con 5 horas de trabajo necesario y 5 de plustrabajo, $\mathit{pl}'=5/5=100~\%$. Si además $c=80$ y $v=20$, $\mathit{pl}=20$, $c/v=4$ y $g'=20/100=20~\%$.
  \item La tasa de ganancia sube con la explotación y baja con la composición orgánica. Es la tasa de beneficio de \autor{Ricardo} \eqref{eq:beneficio-ricardo} cuando todo el capital es variable ($c=0$).
  \end{itemize}
\item \emph{Cómo aumenta la plusvalía}: alargando la jornada (\emph{plusvalía absoluta}) o abaratando los bienes que consume el trabajador gracias a la productividad, que acorta el trabajo necesario (\emph{plusvalía relativa}).
\end{itemize}
\begin{figure}[H]
\centering
\caption{La jornada laboral: trabajo necesario y plustrabajo}\label{fig:jornada}
\grafico[0.7\textwidth]{jornada-laboral}
\fuente{elaboración propia a partir de \textcite[libro~I, caps.~7-9]{Marx1867Kapital}.}
\end{figure}

\subparagraph{iv) Implicaciones}
\begin{itemize}
\item El beneficio, el interés y la renta salen de la plusvalía: son trabajo no pagado. La \concepto{explotación} no es un robo en el intercambio, que es entre equivalentes, sino la consecuencia de que unos tengan los medios de producción y otros solo su fuerza de trabajo.
\item El conflicto entre clases ya estaba en \autor{Ricardo} (salarios frente a beneficios, beneficios frente a rentas); \autor{Marx} lo pone en el centro: cada hora que baja el trabajo necesario es una hora más de plusvalía.
\end{itemize}

\paragraph{c) Acumulación y reproducción}
\begin{itemize}
\item \emph{De la circulación simple al capital}: en la circulación simple se vende para comprar, M--D--M (mercancía, dinero, mercancía): el fin es el valor de uso. El capital invierte el orden, D--M--D$'$: se compra para vender más caro, y el fin es la plusvalía \parencite[libro~I, cap.~4]{Marx1867Kapital}. Por eso la acumulación no tiene límite: empujado por la competencia, el capitalista tiene que acumular (libro~I, cap.~22) \vertema{3.A.34}.
\item \emph{Los esquemas de reproducción} \parencite[libro~II, caps.~20-21]{Marx1885Kapital}: el sector~I produce medios de producción y el~II, bienes de consumo, cada uno con $W_i=c_i+v_i+\mathit{pl}_i$.
  \begin{itemize}
  \item \emph{Reproducción simple}: los capitalistas consumen toda la plusvalía, y el sector~I solo repone los medios de producción de los 2 sectores, $W_{\mathrm I}=c_{\mathrm I}+c_{\mathrm{II}}$, de donde
  \begin{equation}\label{eq:reproduccion-simple}
  v_{\mathrm I}+\mathit{pl}_{\mathrm I}=c_{\mathrm{II}}
  \end{equation}
  lo que los trabajadores y capitalistas del sector~I gastan en consumo es justo lo que el sector~II gasta en medios de producción. En el ejemplo de \autor{Marx}, $c_{\mathrm I}=4.000$, $v_{\mathrm I}=\mathit{pl}_{\mathrm I}=1.000$ y $c_{\mathrm{II}}=2.000$.
  \item \emph{Reproducción ampliada}: parte de la plusvalía se acumula, y el sector~I tiene que producir más de lo que se repone: $v_{\mathrm I}+\mathit{pl}_{\mathrm I}>c_{\mathrm{II}}$.
  \end{itemize}
\item \emph{La lectura}: el crecimiento equilibrado es posible, pero exige una proporción precisa entre los 2 sectores que nadie planea; si se rompe, sobran mercancías en uno y faltan en otro. De estos esquemas parte \autor{Kalecki} para llegar a la demanda efectiva antes que \autor{Keynes} \parencite{Kalecki1933Proba} \vertema{3.A.4}.
\end{itemize}

\paragraph{d) Las leyes del movimiento del capitalismo y las crisis}
\begin{itemize}
\item \autor{Marx} deduce de la acumulación 5 tendencias, que se refuerzan entre sí:
  \begin{itemize}
  \item \emph{Ley de la tendencia decreciente de la tasa de ganancia} \parencite[libro~III, caps.~13-15]{Marx1894Kapital}. Para ganar a sus competidores, cada capitalista mecaniza: sube $c$ frente a $v$. El primero que lo hace gana más; cuando todos lo imitan, la ventaja desaparece, pero la composición orgánica ha subido para todos y, con la tasa de plusvalía dada, $g'=\mathit{pl}'/(1+c/v)$ cae \eqref{eq:tasas-marx}.
    \begin{itemize}
    \item \emph{Contratendencias} (cap.~14): más explotación, capital constante más barato, salarios por debajo de su valor, el ejército de reserva y el comercio exterior. Pero hay un tope: $g'<(v+\mathit{pl})/c$, la tasa que habría con salarios nulos; si el capital crece más deprisa que el trabajo vivo que lo mueve, la tasa de ganancia cae aunque se explote más.
    \item \textcite{Okishio1961Technical} demuestra que si los capitalistas solo adoptan técnicas que les bajan el coste a los precios vigentes y el salario real no cambia, la tasa de ganancia de equilibrio no puede bajar: la ley exige que suban los salarios reales.
    \end{itemize}
  \item \emph{Ley de la concentración y la centralización}: los capitales grandes producen más barato y absorben a los pequeños.
  \item \emph{Ley del ejército industrial de reserva}: la mecanización expulsa obreros, y el paro que crea frena los salarios y da al capital la mano de obra que necesita en las expansiones (libro~I, cap.~23).
  \item \emph{Ley de la miseria creciente}: «la acumulación de riqueza en un polo es al mismo tiempo acumulación de miseria… en el polo opuesto» \parencite[libro~I, cap.~23; traducción propia]{Marx1867Kapital}, en términos relativos o absolutos (los textos admiten las 2 lecturas).
  \item \emph{Ley de las crisis}: cada vez más graves, hasta que el capitalismo cae. «Cría cuervos y te sacarán los ojos»: la burguesía produce, ante todo, sus propios sepultureros \parencite[cap.~I]{MarxEngels1848Manifest}.
  \end{itemize}
\item \emph{Las crisis}, en \autor{Marx}, tienen 3 explicaciones, que no llegó a unir: la \emph{caída de la tasa de ganancia}, que frena la inversión; la \emph{realización}, porque los salarios no bastan para comprar lo que se produce, y la \emph{desproporción} entre los sectores de los esquemas de reproducción. En todas falla la ley de \autor{Say}: con dinero se puede vender sin comprar \vertema{3.A.42}.
\end{itemize}

\paragraph{e) Recomendaciones de política económica}
\begin{itemize}
\item \autor{Marx} apenas describe la sociedad futura. Lo más concreto son las 10 medidas del \emph{Manifiesto} para los países avanzados, entre ellas la expropiación de la tierra, un impuesto fuertemente progresivo, la abolición de la herencia, la centralización del crédito y del transporte en el Estado y la educación pública y gratuita \parencite[cap.~II]{MarxEngels1848Manifest}.
\item El \emph{Manifiesto} termina con «¡Proletarios de todos los países, uníos!», y en la fase superior del comunismo regiría el principio «¡de cada cual, según sus capacidades; a cada cual, según sus necesidades!» \parencite{Marx1875Gotha}.\footnote{Las revoluciones que se reclamaron de \autor{Marx} llegaron a países poco industrializados, sin una teoría de cómo funciona una economía planificada: \autor{Lenin} volvió a un mercado limitado con la Nueva Política Económica (1921); \autor{Stalin} lanzó los planes quinquenales (desde 1928), y \autor{Mao}, el Gran Salto Adelante (1958).}
\end{itemize}

% --------------------------------------------------------------------------
\subsection{Valoración y legado de \autor{Marx}}\label{sec:marx-valoracion}

\begin{itemize}
\item Separando, con términos posteriores (\ref{sec:metodo}), lo positivo de \autor{Marx}, su análisis del capitalismo, de lo normativo, el programa del \emph{Manifiesto}, la valoración empieza por la pieza central del análisis: el paso de los valores a los precios.
\end{itemize}

\paragraph{a) El problema de la transformación}
\begin{itemize}
\item \emph{La pregunta}: si solo el trabajo vivo crea plusvalía, las industrias con mucho capital constante deberían tener una tasa de ganancia menor. Pero la competencia iguala las tasas de ganancia entre industrias. ¿Cómo se concilian?
\item \emph{La respuesta de \autor{Marx}} \parencite[libro~III, cap.~9]{Marx1894Kapital}: las mercancías no se venden por su valor, sino por su \concepto{precio de producción}, el coste ($c+v$) más la ganancia media sobre el capital adelantado. La plusvalía total se reparte entre los capitales según su tamaño, no según el trabajo que emplean, y en el agregado se mantienen 2 igualdades: la suma de los precios es la de los valores, y la de las ganancias, la de las plusvalías (tabla~\ref{tab:transformacion}).
\end{itemize}
\begin{table}[H]
\centering
\caption{De los valores a los precios de producción (tasa de plusvalía del 100~\%)}\label{tab:transformacion}
\small
\begin{tabular}{lcccccccc}
\toprule
\rowcolor{tceeazulpalido}\textbf{Industria} & $c$ & $v$ & $\mathit{pl}$ & $c/v$ & \textbf{Valor} & $g'$ & \textbf{Ganancia} & \textbf{Precio}\\
\midrule
I & 80 & 20 & 20 & 4 & 120 & 20~\% & 30 & \textcolor{tceeazul}{\textbf{130}}\\
II & 70 & 30 & 30 & 2,3 & 130 & 30~\% & 30 & 130\\
III & 60 & 40 & 40 & 1,5 & 140 & 40~\% & 30 & \textcolor{tceerojo}{\textbf{130}}\\
\midrule
Total & 210 & 90 & 90 & & 390 & 30~\% & 90 & 390\\
\bottomrule
\end{tabular}
\fuente{elaboración propia a partir de \textcite[libro~III, cap.~9]{Marx1894Kapital}, que usa 5 industrias. «Ganancia» es la ganancia media (30~\% del capital) y «Precio», el precio de producción, $(c+v)\cdot 1{,}3$; todo el capital constante se consume en el periodo. En azul, el precio por encima del valor (mucho capital constante); en rojo, por debajo.}
\end{table}
\begin{itemize}
\item \emph{Las críticas}:
  \begin{itemize}
  \item \textcite{BohmBawerk1896Abschluss}: la «gran contradicción». El libro~I dice que las mercancías se cambian por su valor; el libro~III, que se cambian por precios distintos. Si los precios no dependen del trabajo, la teoría del valor no explica nada, y con ella cae la explotación \vertema{3.A.3}.
  \item \textcite{Bortkiewicz1907Berichtigung}: \autor{Marx} transformó los precios de los productos, pero dejó en valores los de los medios de producción y los salarios. Hecha bien la transformación, sale un sistema de precios coherente, pero las 2 igualdades no se cumplen a la vez.
  \item \textcite{Samuelson1971Understanding}: para obtener los precios basta con la técnica y el salario real; los valores son un rodeo innecesario.
  \item \textcite{Steedman1977Marx}, con el sistema de \textcite{Sraffa1960Production}: los valores son redundantes, y con producción conjunta puede haber beneficios positivos con plusvalía negativa \parencite{Steedman1975Positive}.
  \end{itemize}
\item \emph{Lo que queda}: \textcite{Morishima1973Marx} demuestra, con \autor{Okishio}, el \emph{teorema marxiano fundamental}: hay ganancias positivas si y solo si la explotación es positiva. Y \textcite{Roemer1982General,Roemer1985Should}, desde el marxismo analítico, propone abandonar el valor trabajo y definir la explotación por la desigualdad en la propiedad de los medios de producción.
\end{itemize}

\paragraph{b) El legado}
\begin{itemize}
\item \emph{En la teoría}: \autor{Kalecki} llega a la demanda efectiva desde los esquemas de reproducción, y \textcite{Robinson1942Essay} lee a \autor{Marx} con \autor{Keynes} y abandona el valor trabajo \vertema{3.A.4}; las teorías de la crisis por la rentabilidad y el subconsumo siguen en el debate sobre el ciclo \vertema{3.A.42}.
\item \emph{Lo que acertó}: la concentración de la industria, la sustitución de trabajo por capital y la recurrencia de las crisis; y la pregunta de cómo se reparte el producto a largo plazo, que vuelve con \textcite{Piketty2014} \vertema{3.A.45}.
\item \emph{Lo que no}: el salario real ha crecido con la productividad en los países industriales, la clase media no desapareció y el capitalismo no se derrumbó; las economías planificadas en su nombre fracasaron en la asignación de recursos.
\item \emph{Recapitulación}: \autor{Marx} cierra la economía política clásica dándole la vuelta: el mismo valor trabajo que en \autor{Ricardo} explicaba el reparto, en \autor{Marx} explica la explotación.
\end{itemize}

% ==========================================================================
\section*{Conclusión}
\minutos{1}

\begin{itemize}
\item \textbf{\emph{Recapitulación}}:
  \begin{itemize}
  \item En conclusión, los clásicos hicieron de la economía una ciencia (\ref{sec:caracterizacion}): con \autor{Smith}, la riqueza nace del trabajo y de su división, y el interés de cada uno, con competencia, la multiplica; con \autor{Ricardo}, el reparto entre salarios, beneficios y rentas decide el crecimiento, que los rendimientos decrecientes de la tierra y la población de \autor{Malthus} empujan hacia el estado estacionario, y que \autor{Mill} ya no ve como una catástrofe (\ref{sec:autores-clasicos}).

  \item \autor{Marx} toma ese aparato, el valor trabajo y el excedente, y le da la vuelta: el beneficio es plusvalía y el capitalismo lleva dentro sus crisis (\ref{sec:marx}).
  \end{itemize}
\item \textbf{\emph{Relevancia}}:
  \begin{itemize}
  \item Sus preguntas siguen abiertas: el reparto de la renta entre el trabajo y el capital \vertema{3.A.45}, el comercio según la ventaja comparativa \vertema{3.B.5} y los límites del crecimiento \vertema{3.B.31}.
  \end{itemize}
\item \textbf{\emph{Extensiones}}:
  \begin{itemize}
  \item La revolución marginalista resuelve la paradoja del agua y los diamantes con la utilidad marginal \vertema{3.A.3}; \autor{Keynes} discute la ley de \autor{Say} y rescata a \autor{Malthus} \vertema{3.A.4}; la \emph{nueva clásica} se llama así porque vuelve a los mercados que se vacían de clásicos y neoclásicos \vertema{3.A.6}, y \autor{Kalecki} lleva a \autor{Marx} al análisis de la demanda efectiva \vertema{3.A.4}.
  \end{itemize}
\item \textbf{\emph{Opinión}}:
  \begin{itemize}
  \item La teoría del valor trabajo fue el talón de Aquiles común: no explicaba los precios observados ni la paradoja del agua y los diamantes, y en \autor{Marx} choca con el problema de la transformación.

  \item Su gran error fue el pesimismo, porque infravaloraron el progreso técnico.

  \item Pero la pregunta de \autor{Ricardo} por las leyes del reparto, que la economía neoclásica convirtió en un problema más de precios, ha vuelto al centro con el debate sobre la desigualdad y la participación del trabajo en la renta.
  \end{itemize}
\item \textbf{\emph{Idea final}}:
  \begin{itemize}
  \item ¿De dónde sale la riqueza y quién se la queda? Del trabajo, de su división y de la acumulación, coinciden todos; en quién se la queda es donde se separan, porque cada uno habla de la feria según le va en ella: \autor{Smith} y \autor{Ricardo}, desde el capitalista que ahorra e invierte; \autor{Malthus}, desde el terrateniente; \autor{Marx}, desde el obrero.
  \end{itemize}
\end{itemize}


\clearpage
\appendix

\section{Anexo: la economía preclásica}\label{anexo:preclasica}

\begin{itemize}
\item El pensamiento anterior a 1776 influyó mucho en la ciencia que nace con \autor{Smith}. Siguiendo a \textcite{LandrethColander2006Historia}, se divide en 2 etapas:
  \begin{itemize}
  \item los \emph{comienzos} (hacia el 800~a.~C.-1500): unos 2.300 años, más de 9 veces los que han pasado desde 1776, con el pensamiento oriental, griego, árabe-islámico y escolástico;
  \item la \emph{era preclásica} (1500-1776): el feudalismo deja paso al comercio, a las ciudades y al Estado nacional; la tierra, el trabajo y el capital empiezan a comprarse en mercados, y la economía empieza a verse como un sistema con sus propias leyes. Son el mercantilismo, los precursores de los clásicos y la fisiocracia.
  \end{itemize}
\end{itemize}

\subsection{Los comienzos: Oriente, Grecia, el islam y los escolásticos}\label{anexo:comienzos}

\begin{itemize}
\item \emph{El pensamiento oriental}: \autor{Schumpeter} no encontró escritos económicos no occidentales con contenido analítico \parencite{Schumpeter1954History}, pero después se han estudiado algunos. En el \emph{Guanzi}, atribuido a \autor{Guan Zhong} (siglo~VII~a.~C.), la «teoría de lo ligero y lo pesado» anticipa la oferta y la demanda (un bien abundante es «ligero» y barato), la teoría cuantitativa (lo mismo aplicado al dinero) y una política de estabilización: que el Estado compre cuando los precios bajan y venda cuando suben.
\item \emph{El pensamiento griego}:
  \begin{itemize}
  \item \autor{Hesíodo} (hacia el siglo~VIII~a.~C.), el primer autor occidental que habla de la escasez y del trabajo;
  \item \autor{Jenofonte} (siglos~V-IV~a.~C.) escribe el \emph{Económico}, de \emph{oikos} (casa) y \emph{nomos} (ley): consejos para gestionar con eficiencia un dominio agrícola, que extiende al ejército y a la administración;
  \item \autor{Platón} (siglos~V-IV~a.~C.) imagina en la \emph{República} y las \emph{Leyes} una ciudad ideal, con una población limitada\footnote{En las \emph{Leyes} fija el número de ciudadanos en 5.040, que es $7!$ y se puede dividir por todos los números del 1 al 12 menos el 11: facilita repartir la tierra y los cargos.} y una moneda que vale por convención, porque la garantiza el poder político;
  \item \autor{Aristóteles} (siglo~IV~a.~C.) distingue el valor de uso del valor de cambio y se pregunta por el precio justo; ve en la moneda algo que resuelve la doble coincidencia de necesidades del trueque; condena el préstamo a interés, porque el dinero no da fruto, y la búsqueda ilimitada de ganancia, y defiende la propiedad privada frente a \autor{Platón}. Sus ideas pasarán a los escolásticos.
  \end{itemize}
\item \emph{El pensamiento árabe-islámico}: los estudiosos árabes no solo transmitieron a \autor{Aristóteles} (los escolásticos lo leyeron traducido del árabe al latín): \autor{al-Ghazali} (siglos~XI-XII) e \autor{Ibn Jaldún} (siglo~XIV), que estudia la tributación y el auge y la caída de las dinastías, analizan la economía de su tiempo dentro de la ley religiosa.
\item \emph{Los escolásticos}: santo \autor{Tomás de Aquino} (siglo~XIII) trae a \autor{Aristóteles} al cristianismo: el \emph{precio justo} es el que fija la estimación común sin fraude ni monopolio, y la usura es ilícita \parencite{Daniel2021Histoire}. Entre los griegos y santo Tomás, \autor{Schumpeter} veía un largo vacío en el análisis económico.
\end{itemize}

\subsection{La Escuela de Salamanca}\label{anexo:salamanca}

\begin{itemize}
\item Teólogos y juristas del siglo~XVI que comentan los contratos de su tiempo y, al hacerlo, analizan los precios y el dinero \parencite{GriceHutchinson1952School}:
  \begin{itemize}
  \item \autor{Francisco de Vitoria} y \autor{Domingo de Soto}, sus fundadores;
  \item \textcite{Azpilcueta1556Cambios} explica la subida de los precios por la llegada de los metales de América: el dinero vale menos donde abunda. Es la teoría cuantitativa, 12 años antes que \textcite{Bodin1568Reponse} \vertema{3.A.35};
  \item \textcite{Mercado1569Tratos} y \textcite{Molina1597Iustitia}: el precio justo es el de mercado, y depende de la abundancia, de la necesidad y del número de compradores;
  \item \autor{Juan de Mariana} critica la alteración de la moneda de vellón como un impuesto encubierto \parencite{Mariana1609Monetae}.
  \end{itemize}
\item No tuvieron continuidad en una ciencia económica propia, porque sus ideas iban dentro de la teología moral; \autor{Schumpeter} los reivindicó entre los primeros que hicieron análisis económico \parencite{Schumpeter1954History}.
\end{itemize}

\subsection{El mercantilismo (1500-1750)}\label{anexo:mercantilismo}

\begin{itemize}
\item \emph{Qué es}: la literatura y la práctica económica de los Estados nacionales en formación. La escriben comerciantes y funcionarios, sobre problemas concretos y a menudo en defensa de un interés; el nombre de «sistema mercantil» se lo ponen sus críticos, entre ellos \autor{Smith} \parencite{LandrethColander2006Historia}.
\item \emph{Ideas}:
  \begin{itemize}
  \item la riqueza total del mundo es fija, y el comercio, un juego de suma cero: un país gana lo que otro pierde;
  \item la riqueza de la nación se mide en metales preciosos y se consigue con una \emph{balanza comercial favorable}: fomentar las exportaciones y frenar las importaciones con aranceles, prohibiciones y subvenciones \parencite{Mun1664England};
  \item salarios bajos para competir y una población numerosa; el poder del Estado como fin.
  \end{itemize}
\item \emph{La evolución}: los primeros mercantilistas, impresionados por el oro y la plata que llegaban de América, querían retener el metal en el país y prohibían exportarlo; los posteriores ven que lo que importa es la balanza en su conjunto: se puede sacar oro para comprar materias primas que luego se exportan transformadas \parencite{LandrethColander2006Historia}.\footnote{\textcite{Chang2002} sostiene que los países hoy ricos se industrializaron con políticas parecidas y luego «retiraron la escalera» al recomendar el libre comercio a los demás \vertema{3.B.28}.}
\item \emph{La crítica}: \autor{Hume} muestra que acumular oro es imposible (el flujo-especie) y \autor{Smith}, que la riqueza es el producto del trabajo (\ref{sec:dinero}). \autor{Keynes} los rehabilita en parte en sus «Notas sobre el mercantilismo», por su preocupación por la escasez de dinero y el tipo de interés \parencite[cap.~23]{Keynes1936}.
\end{itemize}

\subsection{Precursores de los clásicos}\label{anexo:precursores}

\begin{itemize}
\item \autor{William Petty}: la «aritmética política», la medida estadística de la población, la renta y la riqueza de un país \parencite{Petty1690Political}.
\item \autor{John Locke}: el interés y el valor del dinero, en una versión temprana de la teoría cuantitativa \parencite{Locke1692Some}.
\item \autor{Richard Cantillon} (escrito hacia 1730, publicado en 1755): el empresario que compra a precio cierto y vende a precio incierto; el efecto del dinero según por dónde entra en la economía, y una población que sigue a los medios de vida \parencite{Cantillon1755Essai}.
\item \autor{David Hume}: la teoría cuantitativa y el flujo-especie \parencite{Hume1752Money,Hume1752Balance} \vertema{3.B.12}.
\item \autor{James Steuart}, el último gran mercantilista, que defiende la intervención del «estadista» \parencite{Steuart1767Inquiry}.
\item En España, \textcite{Moncada1619Restauracion} atribuye la decadencia a la salida de materias primas y a la entrada de manufacturas extranjeras, y pide protección.
\end{itemize}

\subsection{La fisiocracia}\label{anexo:fisiocracia}

\begin{itemize}
\item \emph{Quién}: \autor{François Quesnay}, médico de la corte de Luis~XV, y sus seguidores (\autor{Mirabeau}, \autor{Dupont de Nemours}). \autor{Turgot}, cercano a ellos, intenta aplicar sus ideas como ministro hasta su caída en 1776 \parencite{HET2026Turgot}.
\item \emph{Ideas}:
  \begin{itemize}
  \item hay un \emph{orden natural} que el gobierno debe respetar; la consigna \emph{laissez faire, laissez passer} se asocia a \autor{Vincent de Gournay}, intendente de comercio próximo a ellos;
  \item solo la agricultura da un \emph{producto neto} (\emph{produit net}), lo que la tierra rinde por encima de lo que cuesta cultivarla; la industria y el comercio solo transforman y mueven lo que la tierra da;
  \item de ahí un \emph{impuesto único} sobre el producto neto, que pagan los propietarios \parencite{HET2026Quesnay}.
  \end{itemize}
\item \emph{El \emph{Tableau économique}} \parencite{Quesnay1758Tableau}: la primera representación del flujo circular de la renta. En el \emph{Analyse} de 1766, \autor{Quesnay} divide la nación en 3 clases \parencite{Quesnay1766Analyse}:
  \begin{itemize}
  \item la \emph{productiva}, los agricultores, que obtienen el producto neto;
  \item la de los \emph{propietarios} de la tierra, que lo cobran como renta;
  \item la \emph{estéril}, artesanos y comerciantes.
  \end{itemize}
  Es el antecedente de los esquemas de reproducción de \autor{Marx} (\ref{sec:marx-aportaciones}) y del equilibrio general \vertema{3.A.21}.
\item \emph{\autor{Turgot}}: explica la formación del capital y el interés \parencite{Turgot1766Reflexions} y formula la ley de las proporciones variables \parencite{Turgot1767Observations} \vertema{3.A.11}.
\item \emph{Influencia}: \autor{Smith} los trata en París y toma de ellos el \emph{laissez-faire}, el excedente (que extiende a la industria) y la división en 3 clases.
\end{itemize}

\subsection{El paso a la era clásica}\label{anexo:paso}

\begin{itemize}
\item 1776 es una fecha convencional, y con mirada británica: los escolásticos de Salamanca y los fisiócratas ya habían estudiado la economía con método, pero no dejaron una disciplina con continuidad. \autor{Smith} reúne sus piezas en un sistema y funda una escuela.
\end{itemize}

\section{Anexo: los clásicos y \autor{Marx} por líneas de investigación}\label{anexo:lineas}

\begin{itemize}
\item La ficha del tribunal propone ordenar el tema por líneas de investigación. Sin cambiar la estructura de la exposición, la tabla~\ref{tab:lineas} da esa vista cruzada, útil para el repaso y para las preguntas del tribunal.
\end{itemize}
\begin{table}[H]
\centering
\caption{Las líneas de investigación, autor por autor}\label{tab:lineas}
\scriptsize
\begin{tabular}{>{\raggedright\arraybackslash}m{1.15cm}>{\raggedright\arraybackslash}m{1.85cm}>{\raggedright\arraybackslash}m{2.1cm}>{\raggedright\arraybackslash}m{1.65cm}>{\raggedright\arraybackslash}m{1.85cm}>{\raggedright\arraybackslash}m{1.75cm}>{\raggedright\arraybackslash}m{1.85cm}}
\toprule
\rowcolor{tceeazulpalido} & \textbf{Valor} & \textbf{Crecimiento y distribución} & \textbf{Ley de \autor{Say}} & \textbf{Dinero} & \textbf{Comercio} & \textbf{Política}\\
\midrule
\autor{Smith} & Trabajo comandado; precio natural como suma de costes & División del trabajo; fondo de salarios; la renta, residuo & Lo ahorrado se gasta & Letras reales & Ventaja absoluta & Los 3 deberes del soberano; las 4 máximas\\
\addlinespace[0.3em]
\autor{Say} & Utilidad & El empresario & La formula & Medio de cambio & Libre cambio & \emph{Laissez-faire}\\
\addlinespace[0.3em]
\autor{Ricardo} & Trabajo incorporado en la tierra peor & Renta diferencial; beneficio residual; estado estacionario & La acepta & Bullionista; patrón en lingotes & Ventaja comparativa & Contra las Leyes de Granos\\
\addlinespace[0.3em]
\autor{Malthus} & Oferta y demanda & Población; subsistencia & La niega: exceso general & --- & Leyes de Granos & Consumo improductivo\\
\addlinespace[0.3em]
\autor{Mill} & Los 3 casos según la oferta & Reparto institucional; estado estacionario deseable; fin del fondo de salarios & Igualdad, no identidad & Flujo-especie con tipos; \emph{banking school} & Demanda recíproca; industria naciente & Reformas; mínimo exento\\
\addlinespace[0.3em]
\textcolor{tceerojo}{\autor{Marx}} & Trabajo socialmente necesario; $W=c+v+\mathit{pl}$ & Plusvalía; tasa de ganancia decreciente & La niega: crisis & D--M--D$'$ & --- & Las 10 medidas del \emph{Manifiesto}\\
\bottomrule
\end{tabular}
\fuente{elaboración propia a partir de los apartados~\ref{sec:autores-clasicos} y~\ref{sec:marx}.}
\end{table}

\section{Anexo: la ventaja comparativa de \autor{Ricardo} en extenso}\label{anexo:ricardo-comercio}

\begin{itemize}
\item Este anexo desarrolla, en clave de historia del pensamiento, lo resumido en el apartado~\ref{sec:ricardo-comercio}. El modelo de referencia, con sus gráficos y extensiones, es del tema del comercio internacional \vertema{3.B.5}.
\end{itemize}

\paragraph{a) Idea}
\begin{itemize}
\item Con la ventaja absoluta de \autor{Smith}, un país menos productivo en todo quedaría fuera del comercio. \autor{Ricardo} muestra que el comercio se basa en las diferencias \emph{relativas} de costes, que reflejan diferencias de tecnología: cada país exporta el bien que produce con menor coste relativo, es decir, con menor coste de oportunidad \parencite[cap.~7]{Ricardo1821Principles}.
\item El contexto es el de las Leyes de Granos (\ref{sec:ricardo-contexto}): la ventaja comparativa es un argumento para importar grano.
\end{itemize}

\paragraph{b) Modelo}
\subparagraph{i) Supuestos}
\begin{itemize}
\item 2 países (A, Inglaterra; B, Portugal), 2 bienes homogéneos (paño, $X$; vino, $Y$) y 1 factor, el trabajo, homogéneo y divisible, móvil dentro de cada país e inmóvil entre países.
\item Tecnologías distintas entre países y lineales: los costes unitarios en trabajo, $a_X^A$, $a_Y^A$, $a_X^B$, $a_Y^B$, son constantes (rendimientos constantes a escala; producto marginal igual al medio).
\item Teoría del valor trabajo: dentro de cada país, el precio relativo es el cociente de los costes en trabajo.
\item Libre comercio, sin costes de transporte.
\end{itemize}

\subparagraph{ii) Desarrollo}
\begin{itemize}
\item \emph{Condición necesaria}: que los costes relativos difieran, $a_X^A/a_Y^A\neq a_X^B/a_Y^B$. Con las cifras de la tabla~\ref{tab:pano-vino}, $100/120=0{,}83<90/80=1{,}125$: Inglaterra tiene ventaja comparativa en paño y Portugal, en vino (Portugal tiene ventaja absoluta en los 2, pero mayor en vino: $80/120<90/100$).
\item \emph{Condición suficiente}: que la relación real de intercambio quede \emph{estrictamente} entre los 2 costes relativos \eqref{eq:ventaja-comparativa}. Si coincide con uno de ellos, ese país obtiene con el comercio lo mismo que produciendo él mismo y no gana nada; si queda fuera del intervalo, uno de los 2 pierde.
\end{itemize}

\subparagraph{iii) Implicaciones}
\begin{itemize}
\item \emph{Por qué se comercia}: por las diferencias de precios relativos que causan las diferencias de tecnología.
\item \emph{Qué se produce}: cada país solo el bien de su ventaja comparativa (especialización completa, por la tecnología lineal).
\item \emph{Bienestar}: ganan los 2, porque la producción total aumenta y cada uno consume por encima de lo que podría producir solo; el comercio amplía sus posibilidades de consumo.
\item \emph{La relación real de intercambio}: el modelo, solo de oferta, no la determina.
\end{itemize}

\paragraph{c) Valoración}
\begin{itemize}
\item \emph{El mérito}: dar las condiciones necesaria y suficiente del comercio beneficioso. Que no determine la RRI importa poco si se lee, como propone \textcite{Bhagwati1964Pure}, como una demostración de que el comercio es beneficioso.
\item \emph{La demanda}: la cierra \textcite{Mill1844Essays} con la demanda recíproca (\ref{sec:mill-comercio}), y \textcite{Marshall1879Pure} la dibuja con las curvas de oferta recíproca, que llevan a la teoría neoclásica del comercio.
\item \emph{La tecnología lineal}: con 2 factores y funciones de producción neoclásicas, la especialización deja de ser completa \vertema{3.B.5}.
\item \emph{La autoría}: un folleto anónimo de 1814 y \textcite{Torrens1815Essay} se adelantan; la discusión sigue abierta \parencite{GrancayGrancay2015Considerations,Ruffin2002Ricardo}.
\end{itemize}

\section{Anexo: aspectos para el test}\label{anexo:test}

\begin{itemize}
\item El banco del simulador tiene 12 preguntas oficiales de este tema. Lo que preguntan y dónde se trata:
\end{itemize}
\begin{table}[H]
\centering
\caption{Las preguntas del test sobre los clásicos y \autor{Marx}}\label{tab:test}
\scriptsize
\begin{tabular}{>{\raggedright\arraybackslash}m{2.2cm}>{\raggedright\arraybackslash}m{9.4cm}>{\raggedright\arraybackslash}m{2.4cm}}
\toprule
\rowcolor{tceeazulpalido}\textbf{Examen} & \textbf{Idea que se pregunta (respuesta)} & \textbf{Dónde}\\
\midrule
2002 (octubre) & Con 3 tierras que necesitan 10, 15 y 20 horas por unidad, el trigo vale 20: lo fija la tierra peor que se cultiva & \ref{sec:ricardo-valor}\\
2006 (septiembre) & Los clásicos usaron el valor trabajo para buscar una medida invariable del valor; no lo inventa \autor{Marx} & \ref{sec:ricardo-valor}\\
2009 (octubre) & En la \emph{Riqueza}: crecimiento por la división del trabajo, limitada por el mercado y el capital, y supresión de trabas (todas ciertas) & \ref{sec:smith-riqueza}\\
2013 (noviembre) & El comercio de \autor{Mill} aporta ideas que desarrollará \autor{Marshall} & \ref{sec:mill-comercio}\\
2014 (septiembre) & Para el tribunal, el valor en \autor{Smith} se basa en el trabajo incorporado (en rigor, solo en el estado primitivo; la medida es el trabajo comandado) & \ref{sec:smith-valor}\\
2016 (enero) & Para \autor{Mill}, el reparto de las ganancias del comercio depende de la intensidad relativa de las demandas & \ref{sec:mill-comercio}\\
2016 (noviembre) & En \autor{Marx}, valores y precios difieren por las diferencias en la composición orgánica del capital & \ref{sec:marx-valoracion}\\
2018 (enero) & Es \emph{falso} que \autor{Mill} quisiera una economía más normativa que positiva & \ref{sec:mill}, \ref{sec:metodo}\\
2020 (enero) & Los clásicos atendieron sobre todo a la interacción entre la distribución de la renta y el producto & \ref{sec:schumpeter}\\
2022 (julio) & Es \emph{falso} que para \autor{Smith} y \autor{Ricardo} el beneficio sea el residuo: en \autor{Smith} lo es la renta & \ref{sec:smith-distribucion}\\
2024 (junio) & Es \emph{falso} que \autor{Ricardo} fuera contrario al bullionismo: fue el gran bullionista & \ref{sec:ricardo-say}, \ref{sec:dinero}\\
2025 (marzo) & Es \emph{falso} que el precio natural de \autor{Smith} sea el de la oferta y la demanda y el de mercado el de los costes: es al revés & \ref{sec:smith-valor}\\
\bottomrule
\end{tabular}
\fuente{banco de preguntas del simulador de test.}
\end{table}
\begin{itemize}
\item \emph{Otras trampas}: la mano invisible sale una sola vez en la \emph{Riqueza}; el nombre «ley de hierro de los salarios» es de \autor{Lassalle} (1863); el de «clásicos», de \autor{Marx}; \autor{Ricardo} murió en 1823, antes del debate entre la \emph{currency} y la \emph{banking school}; \autor{Keynes} no dedicó la \emph{Teoría general} a \autor{Malthus}.
\end{itemize}

\begin{anotaciones}
Preguntas tipo test del primer ejercicio: \href{https://www.victorgutierrezmarcos.es/oposicion/temario/primer-ejercicio/test/simulador.html}{simulador de test}.

\textbf{Para ampliar}:
\begin{itemize}
\item The History of Economic Thought Website: «The Ricardian Era» y «The Bullionist Controversies» \parencite{HET2026Ricardian,HET2026Bullion}.
\item \autor{Eduardo Escartín}, apuntes sobre \autor{Ricardo} (Universidad de Sevilla): \url{https://personal.us.es/escartin/Ricardo.pdf} \parencite{Escartin2008Apuntes}.
\item Documentales sobre \autor{Adam Smith} (Free To Choose Network): \url{https://www.youtube.com/watch?v=V6S6pMsKzlI} y \url{https://www.youtube.com/watch?v=8ruiUOQERnw}.
\end{itemize}
\end{anotaciones}

\bibliografia
\end{document}
